%---------------------------------------------------------------------------%
%-                                                                         -%
%-                           LaTeX Template                                -%
%-                                                                         -%
%---------------------------------------------------------------------------%
% 我加入的
%- Copyright (C) Huangrui Mo <huangrui.mo@gmail.com>
%- This is free software: you can redistribute it and/or modify it
%- under the terms of the GNU General Public License as published by
%- the Free Software Foundation, either version 3 of the License, or
%- (at your option) any later version.
%---------------------------------------------------------------------------%
%->> Document class declaration
%---------------------------------------------------------------------------%
\documentclass[printcopy,fontset=windows]{Style/neuthesis}%
%- Multiple optional arguments:
%- [<singlesided|doublesided|printcopy>]% set one or two sided eprint or print
%- [draftversion]% show draft version information
%- [fontset=<fandol|windows|mac|adobe>]% specify font set to replace automatic detection
%- [scheme=plain]% thesis writing of international students
%- [standard options for ctex book class: draft|paper size|font size|...]%
%---------------------------------------------------------------------------%
%->> Document settings
%---------------------------------------------------------------------------%
\usepackage[most]{tcolorbox}
\usepackage[bibtex,myhdr,table,tikz,list,geometry,math]{Style/artratex}
% document settings
%- usage: \usepackage[option1,option2,...,optionN]{artratex}
%- Multiple optional arguments:
%- [bibtex|biber]% set bibliography processor and package
%- [geometry]% reconfigure page layout via geometry package
%- [lscape]% provide landscape layout environment
%- [myhdr]% enable header and footer via fancyhdr package
%- [color]% provide color support via xcolor package
%- [background]% enable page background
%- [tikz]% provide complex diagrams via tikz package
%- [table]% provide complex tables via ctable package
%- [list]% provide enhanced list environments for algorithm and coding
%- [math]% enable some extra math packages
\usepackage{Style/artracom}% user defined commands
% 我加入的
\usepackage{algorithm}
% ====== 占位符宏（方便学弟搜索 "TBD" 定位所有待填内容） ======
% 用法：\TBD{期望填入的描述} ；编译后显示为红色粗体，醒目提示。
% 所有数字/图/表的待填处均已用此宏标注。
\providecommand{\TBD}[1]{\textcolor{red}{\textbf{[待填:#1]}}}
\providecommand{\TBDFIG}[1]{%
  \fbox{\parbox{0.9\linewidth}{\centering\vspace{1em}%
  \textcolor{red}{\textbf{[待填图:#1]}}\vspace{1em}}}%
}
% ==========================================================
\usepackage{algorithmicx}
\usepackage{algpseudocode}
\usepackage{arydshln}  % 支持虚线
% \usepackage{subfigure}  % 必须加上！


%这个图用到了\usepackage{pgfplots} 和 \usetikzlibrary{patterns}这两个包

%---------------------------------------------------------------------------%
%->> Document inclusion
%---------------------------------------------------------------------------%
%\includeonly{Tex/Chap_1,...,Tex/Chap_N}% selected files compilation
%---------------------------------------------------------------------------%
%->> Document content
%---------------------------------------------------------------------------%
\def\alltex{}
% bibtex cache-buster: bumping this tag forces Overleaf to re-run bibtex on
% the next compile so new entries in Biblio/ref.bib are picked up.
% buster: 2026-04-25-02
\usepackage{tikz}       % TikZ 是基礎
\usepackage{pgfplots}
\pgfplotsset{compat=1.18}% silence "backwards compatibility mode" warning
\begin{document}

% Force every citation used in the query-rewriting chapters to be picked up
% by bibtex even when an incremental build reuses a stale output.bbl.  All
% keys below are defined in Biblio/ref.bib.
\nocite{mo2024convsar,lee2017coref,halliday1976cohesion,quan2019gecor,
  elgohary2019canyouunpack,vakulenko2021question,mao2023llm4cs,
  voskarides2020quretec,yu2020fewshotgenerative,wu2022conqrr,
  mo2023convgqr,jang2024itercqr,lin2020conversational,raffel2020t5,
  ye2023enhancing,mao2023search,kim2022saving,schulman2017ppo,
  lin2023dragon,warner2024modernbert,qu2020orquac,mtrsuite2025,lora2022}
\begin{sloppypar}


{
\titlecontents{chapter}
    [1.55cm]
    {\vspace{-1.2ex}\bf\large}
    {\contentslabel{3em}}
    {\hspace*{-3em}} % 无编号项目格式
    % {{\titlerule*[0.3pc]{\raisebox{0.5ex}{.}}\normalfont\chapterPageNumber\contentspage}}
    {% 创建一个局部作用域以应用我们的无加粗点线定义
  \mytitlerule
  \chapterPageNumber
  \contentspage % 页码
  }
\titlecontents{section}
    [1.4cm]
    {\vspace{-1.2ex}\songti\normalsize}
    {\contentslabel{2.5em}}{}
    {\titlerule*[0.3pc]{\raisebox{0.5ex}{.}}\contentspage}

\titlecontents{subsection}
    [2.6cm]{\vspace{-1.2ex}\songti}
    {\contentslabel{2.5em}}{}
    {\titlerule*[0.3pc]{\raisebox{0.5ex}{.}}\contentspage}

}


% \layout
%-
%-> Frontmatter: title page, abstract, content list, symbol list, preface
%-
\frontmatter% initialize the environment
% !TEX ROOT = ../Thesis.tex 
%---------------------------------------------------------------------------%
%->> 封面信息及生成
%---------------------------------------------------------------------------%
%-
%-> 中文封面信息
%-
\category{}%分类号
\confidential{}% 密级
\UDC{}
\schoollogo% 校徽

\title{面向多轮对话检索的查询重写方法研究}% 论文中文题目
\author{}% 论文作者
\advisor{}% 指导教师
\advisorsec{计算机科学与工程学院}%
\degree{硕士}%
\major{计算机技术}%
\institute{计算机科学与工程学院}%

\chinesetitle{}
\research{}%
\authorno{}%
\chinesedate{2026~年~6~月}%
\submissiondate{2026年4月}%
\oraldefencedate{2026年5月}%
\degreedate{2026年6月}%
\chairman{}%
%-
%-> 英文封面信息
%-
\englishtitle{Research on Query Rewriting Methods for Multi-turn Conversational Retrieval}
\englishauthor{}%
\englishadvisor{}%
\englishdegree{Master}%
\englishdegreetype{Engineering}%
\englishthesistype{thesis}%
\englishuniversity{Northeastern University}
\englishmajor{Computer Technology}%
\englishinstitute{School of Computer Science and Engineering}%
\englishdate{June 2026}%
%-
%-> 生成封面
%-
\makeprintcover
\makechinesetitle
\makeenglishtitle
%-
%-> 作者声明
%-
\makedeclaration
% title page, abstract, dedication
% !TEX ROOT = ../Thesis.tex 
%-
%-> 中文摘要
%-
\chapter{摘\quad 要}\chaptermark{摘\quad 要}
\linespread{1.5}
\zihao{-4}

在大语言模型与检索增强生成（Retrieval-Augmented Generation, RAG）技术迅猛发展的背景下，多轮对话检索已成为对话式人工智能系统落地应用的关键环节。然而，真实多轮对话中普遍存在的指代、省略、话题漂移等语言现象，使得当前轮次的查询往往语义残缺，无法被现有检索器直接理解和召回。查询重写（Query Rewriting）作为在检索之前将上下文依赖的用户问句改写为语义自包含独立查询的关键预处理模块，直接决定了下游检索与生成的上限。然而，现有查询重写方法仍面临三大挑战：其一，对长历史、跨话题对话中的噪声轮次缺乏显式建模，导致重写时引入大量无关信息；其二，指代消解与省略补全被混合为端到端序列生成任务，训练目标不明确且错误难以定位；其三，重写目标局限于"像人写的参考重写"，与下游检索器的真实偏好存在系统性偏差。针对上述痛点，本文系统研究了面向多轮对话检索的查询重写方法，提出从"历史剪枝—渐进式重写—检索反馈优化"三阶段构建查询重写框架，并在多个公开基准上进行了深入的实证分析。本文的主要研究工作与创新成果如下：

首先，针对多轮对话中历史轮次的话题漂移问题，提出了一种话题漂移感知的对话历史剪枝方法。不同于将完整历史一视同仁地拼接输入的传统做法，本方法在轮次级引入三分类相关性判别器（强相关 / 弱相关 / 无关），基于 ModernBERT-base 构建轻量级分类头，并利用 TopiOCQA 的话题段落标注与 QReCC 的人工重写作为弱监督训练信号。在 94\,422 条轮次-查询对上训练后，分类器在留出验证集上取得 80.1\% 的准确率与 0.797 的宏 F1，能够稳定地压缩历史、降低输入噪声。

其次，针对指代消解与省略补全耦合导致错误难以分离的问题，提出了一种两阶段渐进式查询重写框架（Progressive Rewriter, PRW）。该框架将重写解耦为"指代消解"与"省略补全"两个串行子任务：第一阶段仅将指代词替换为对应实体，保持表层结构不变；第二阶段在此基础上补全缺失的主语、谓语与修饰成分，生成自包含查询。基于 Qwen2.5-3B-Instruct 与 LoRA（$r{=}16$）的两阶段 SFT 实验表明，PRW 在 TopiOCQA、Doc2Dial 两个话题切换密集、上下文依赖强的基准上，将 Recall@5 相对 No-Rewrite 分别提升 39 和 66 个百分点以上（BGE-large 检索器下）。

最后，针对重写目标与下游检索器偏好失配的问题，提出了检索反馈驱动的强化微调方法（Retrieval-Feedback Rewriter, RFR）。本文从 TopiOCQA 训练集的 10\,000 个带 gold passage 的轮次采样 80\,000 条候选重写，以 BGE-large 计算"候选 vs gold passage"余弦相似度作为反馈，构造 4\,106 条高质量偏好对，再用 Direct Preference Optimization (DPO) 在 PRW 基础上继续 LoRA 训练。在 129 步训练内，reward accuracies 从 45\% 单调上升至 85\%，训练全程稳定。在 QReCC、TopiOCQA、Doc2Dial、QuAC 四个公开基准 $\times$ BGE-large / GTE-ModernBERT / Dragon-Multiturn 三种异构检索器共 12 组设定中，RFR 相对 PRW 取得平均 $+1.10$ R@5 / $+1.31$ R@20 的一致增益，且该增益在训练从未见过的 GTE-ModernBERT 与 Dragon-Multiturn 上同样成立。

综上所述，本文围绕多轮对话检索中的查询重写问题，构建了一套"噪声过滤—结构化重写—下游反馈"相互衔接的技术方案，深入研究了对话历史建模、渐进式生成与基于检索反馈的偏好优化等核心问题，为工业级对话式 RAG 系统提供了一种轻量、高效且鲁棒的查询重写范式。

\keywords{多轮对话检索；查询重写；检索增强生成；直接偏好优化；大语言模型}

%-
%-> 英文摘要
%-
\chapter{Abstract}\chaptermark{Abstract}

Against the backdrop of the rapid progress of large language models (LLMs) and retrieval-augmented generation (RAG), multi-turn conversational retrieval has become a critical building block for deploying conversational AI systems. In realistic multi-turn dialogues, however, linguistic phenomena such as coreference, ellipsis, and topic drift are ubiquitous, rendering the current-turn query semantically incomplete and hard for existing retrievers to understand. Query rewriting (QR), which transforms a context-dependent utterance into a self-contained standalone query before retrieval, directly determines the upper bound of downstream retrieval and generation. Nevertheless, existing QR approaches still suffer from three key limitations: (i) they lack explicit modeling of noisy turns in long, topic-switching histories, so that irrelevant information is easily leaked into the rewrite; (ii) coreference resolution and ellipsis completion are conflated into a single end-to-end sequence generation task with ill-defined training signals; and (iii) the training objective is tied to "human-written references" and is systematically misaligned with the true preferences of downstream retrievers. To address these issues, this thesis systematically investigates query rewriting methods for multi-turn conversational retrieval, and proposes a three-stage framework consisting of \emph{history pruning}, \emph{progressive rewriting}, and \emph{retrieval-feedback optimization}. The main contributions are as follows.

First, to tackle topic drift in long conversational histories, we propose a \textbf{topic-drift-aware history pruning} module. The method introduces a turn-level three-way relevance classifier (strongly-, weakly-, and un-related) built on top of ModernBERT-base, trained with weak supervision derived from the topic-section labels of TopiOCQA and the human-written rewrites of QReCC. Trained on 94{,}422 turn--query pairs, the classifier achieves 80.1\% accuracy and macro-F1 = 0.797 on a held-out validation set, reliably compressing history and reducing input noise.

Second, to resolve the coupling between coreference and ellipsis phenomena, we propose a \textbf{Progressive Rewriter (PRW)} framework that decomposes rewriting into two sequential subtasks: a coreference-only stage that replaces anaphoric mentions while preserving surface structure, and an ellipsis-completion stage that restores missing subjects, predicates, and modifiers. With Qwen2.5-3B-Instruct fine-tuned via LoRA ($r{=}16$), PRW raises Recall@5 by more than 39 and 66 percentage points over No-Rewrite on TopiOCQA and Doc2Dial respectively (BGE-large retriever), i.e., the two benchmarks where topic switches and context dependencies are most severe.

Third, to bridge the gap between human-style rewrites and retriever-friendly rewrites, we design a \textbf{Retrieval-Feedback Rewriter (RFR)}. We sample 80{,}000 candidate rewrites from 10{,}000 TopiOCQA training turns with gold passages, score each candidate by its BGE-large cosine similarity to the gold passage, and assemble 4{,}106 high-quality preference pairs. DPO is then applied on top of PRW's LoRA adapter. Within only 129 optimization steps, reward accuracies rise monotonically from 45\% to 85\%. On 12 settings (four public benchmarks $\times$ three heterogeneous retrievers), RFR delivers a consistent average gain of $+1.10$ Recall@5 / $+1.31$ Recall@20 over PRW, and the gain transfers to unseen retrievers (GTE-ModernBERT and Dragon-Multiturn).

In summary, this thesis presents a unified technical route that couples noise filtering, structured rewriting, and downstream feedback, and provides a lightweight, efficient and robust query-rewriting paradigm for industrial-scale conversational RAG systems.

\englishkeywords{multi-turn conversational retrieval; query rewriting; retrieval-augmented generation; direct preference optimization; large language models}
 % abstract (chinese and english)

{% content list region
% \linespread{1}
% \intotoc{\contentsname}% add link to contents table and bookmark
% \chapter*{\contentsname}
\linespread{1.5}% local line space
\tableofcontents% contents catalog
{
\titlecontents{chapter}
    [1.55cm]        % 调整章标题左侧缩进值
    {\vspace{-1.2ex}\bf\large}
    {\contentslabel{3em}}
    {\hspace*{-3em}} % 无编号项目格式
    % {{\titlerule*[0.3pc]{\raisebox{0.5ex}{.}}\normalfont\chapterPageNumber\contentspage}}
    {% 创建一个局部作用域以应用我们的无加粗点线定义
  \mytitlerule
  \chapterPageNumber
  \contentspage % 页码
  }
\titlecontents{section}
    [1.4cm]
    {\vspace{-0.8ex}\songti\normalsize}
    {\contentslabel{1.5em}}{}
    {\titlerule*[0.3pc]{\raisebox{0.5ex}{.}}\contentspage}

\titlecontents{subsection}
    [2.6cm]{\vspace{-1.2ex}\songti}
    {\contentslabel{2.5em}}{}
    {\titlerule*[0.3pc]{\raisebox{0.5ex}{.}}\contentspage}

}
% \intotoc{\listfigurename}% add link to contents table and bookmark
% \listoffigures% figures catalog

% \intotoc{\listtablename}% add link to contents table and bookmark
% \listoftables% tables catalog
}
% \input{Tex/Prematter}% list of symbols, preface content
%-
%-> Mainmatter
%-
\mainmatter% initialize the environment
\linespread{1.4}
\zihao{-4}


%---------------------------------------------------------------------------%
%->> Main content
%---------------------------------------------------------------------------%
\chapter{绪论}
\label{chap:introduction}

\section{研究背景与意义}

\subsection{检索增强生成的发展现状}

大语言模型（Large Language Models, LLMs）的快速发展是近年来人工智能领域最引人注目的技术突破之一。以 GPT-4~\citep{openai2023gpt4}、Gemini~\citep{geminiteam2024gemini15}、DeepSeek-V3~\citep{deepseekai2024deepseekv3}、Qwen2.5~\citep{qwen2025qwen25technicalreport} 为代表的大语言模型，在文本理解、逻辑推理、代码生成等广泛的自然语言处理任务上展现出了令人瞩目的能力，深刻改变了人机交互的范式。然而，尽管大语言模型的规模不断扩大、能力持续增强，其固有的局限性也日益凸显。

首先，知识时效性受限。大语言模型的知识来源于预训练阶段所接触的静态语料库，存在明确的"知识截止日期"，无法获取训练数据截止日期之后的新信息和新事件。其次，事实性幻觉问题严重。研究表明，大语言模型在生成过程中倾向于产生看似流畅合理但实际上与客观事实不符的内容~\citep{xu2025hallucination}，这种现象在医疗、法律、金融等对准确性要求极高的领域构成了严重的安全风险。第三，长尾知识与领域知识匮乏。Kandpal 等人的研究指出，大语言模型在学习长尾知识方面存在显著困难~\citep{kandpal2023largelanguagemodelsstruggle}——对于训练语料中出现频率较低的实体和事实，模型的记忆与生成能力急剧下降。此外，对于特定垂直领域或企业私有知识库中的专业知识，通用大语言模型往往缺乏足够的覆盖。

为了克服上述局限性，检索增强生成（Retrieval-Augmented Generation, RAG）技术应运而生，并迅速成为当前大语言模型应用中最核心的技术范式之一~\citep{gao2024retrievalaugmented}。RAG 的核心思想是在模型生成回答之前，首先从外部知识库中检索出与用户查询相关的文档或段落，然后将这些检索到的证据作为上下文提供给语言模型，从而引导其生成更加准确、有据可查的回复。这种将参数化记忆（模型内部权重）与非参数化记忆（外部知识库）相结合的思路，使得模型可以动态地访问最新信息，有效缓解了知识过时和幻觉问题~\citep{shuster2021retrieval}。Lewis 等人~\citep{lewis2020rag} 与 Guu 等人~\citep{guu2020realm} 的里程碑式工作奠定了现代 RAG 研究的基本范式，此后 RAG 已从辅助手段成长为连接大模型能力与行业知识的核心基础设施。

从技术架构来看，现代 RAG 系统已经从最初的"检索-拼接-生成"朴素流水线，演化为包含查询理解、多路检索、重排序、上下文压缩、自适应检索等多个精细化模块的复杂系统。其中，\textbf{检索模块的召回能力决定了整个系统的上限}——即使生成端拥有再强的语言模型，若检索到的证据不相关或不完整，后续生成无法无中生有。以 DPR~\citep{karpukhin-etal-2020-dense}、Contriever~\citep{izacard2022contriever}、BGE~\citep{xiao2023cpack} 为代表的密集向量检索模型取代了传统的稀疏检索（如 BM25），在单轮场景下已具备高度成熟的语义匹配能力。

\subsection{多轮对话检索与查询重写的核心挑战}

随着 ChatGPT 等对话式产品的普及，用户与 RAG 系统的交互方式已从单轮独立查询向\textbf{多轮对话}形态快速迁移。在真实场景下，用户往往不会一次性把所有信息表达清楚，而是通过连续多轮追问、澄清、切换话题来逐步完成信息获取。这种\textit{多轮对话检索}（Multi-turn Conversational Retrieval）场景给系统带来了远比单轮更为复杂的挑战。

考察以下一个真实多轮对话片段：

\begin{quote}
\textbf{U1}：What is the capital of France? \\
\textbf{A1}：Paris. \\
\textbf{U2}：Who founded it? \\
\textbf{A2}：Paris was founded by the Parisii tribe in the 3rd century BC. \\
\textbf{U3}：And its population now?
\end{quote}

当检索系统需要为 U3 检索相关证据时，直接把 "And its population now?" 输入检索器将无法得到任何有意义的结果——"its" 指代不明，句子的主语被省略，信息完全依赖对话历史。只有把查询补全为"What is the current population of Paris?" 这样的自包含查询，现代稠密检索器才能返回正确的证据文档。这种\textbf{把上下文依赖的用户问句改写为语义自包含的独立查询}的技术，即为\textbf{查询重写}（Query Rewriting, QR），是多轮对话 RAG 系统中承上启下的关键环节~\citep{mo2024convsar}。

真实多轮对话中带来查询语义不完整的语言现象主要包括：
\begin{itemize}
\item \textbf{指代（Coreference）}：使用"他/她/它/前者/这个"等代词替代先前提到的实体~\citep{lee2017coref}；
\item \textbf{省略（Ellipsis）}：省略主语、谓语、宾语或修饰成分，依赖上下文补全~\citep{halliday1976cohesion,quan2019gecor}；
\item \textbf{话题漂移（Topic Drift）}：后续轮次可能与早期轮次话题不再直接相关，甚至发生硬切换~\citep{adlakha-etal-2022-topiocqa}；
\item \textbf{隐式澄清}：用户基于系统前一轮回答进行追问，查询的含义依赖于答案的内容。
\end{itemize}

针对上述挑战，业界主要存在两条技术路线。一条是\textit{端到端对话检索器}（如 Dragon-ChatQA~\citep{Liu2024ChatQASG}），直接把对话历史与当前轮次拼接作为检索输入，让模型自行学会消解上下文依赖；另一条则是\textit{查询重写}~\citep{elgohary2019canyouunpack,vakulenko2021question,mao2023llm4cs}，在检索之前显式地把当前轮次改写为独立查询，再交给任意现成的单轮检索器。本文聚焦于后者，原因在于：（1）查询重写与下游检索器完全解耦，可即插即用地与任意单轮检索器组合，部署灵活性显著高于专用对话检索器；（2）重写后的查询可直接复用于排序、摘要、用户意图分析等多个下游模块，具有更强的工程复用价值；（3）现有对话检索器的训练依赖昂贵的对话标注数据，而查询重写方法可以充分利用日益丰富的大语言模型上下文理解能力，研究空间更大。

\subsection{现有查询重写方法的瓶颈}

尽管围绕查询重写已有一批代表性工作~\citep{voskarides2020quretec,yu2020fewshotgenerative,wu2022conqrr,mao2023llm4cs,mo2023convgqr,jang2024itercqr}，现有方法仍远未解决多轮对话检索的核心难点。本文将其归纳为三大瓶颈：

\textbf{瓶颈一：历史噪声建模不足。}绝大多数已有方法在重写阶段将对话历史 $H_i=\{(q_1,a_1),\dots,(q_{i-1},a_{i-1})\}$ 作为完整上下文输入重写模型。但在长对话、跨话题场景下，远早于当前轮次的历史往往已与当前查询无关，甚至存在硬话题切换——此时继续引入无关历史不仅增大输入长度、推高计算成本，更会系统性误导重写器引入无关实体或属性，造成"重写越长、检索越差"的反直觉现象。少数工作尝试用手工规则或启发式截断，但无法区分"同一话题内的弱相关轮次"与"已切换到新话题的噪声轮次"，缺乏对话题漂移的显式建模。

\textbf{瓶颈二：异质现象耦合训练。}现有方法普遍把指代消解与省略补全混合为一个端到端序列生成任务，让模型在单一损失下同时学习两种异质的语言现象。然而，这两种现象在语言学性质、对检索的影响、错误传播模式上差异显著——指代错误通常表现为错词（wrong entity），省略错误通常表现为缺词（missing slot）。端到端训练难以让模型明确地学习两种结构不同的修复操作，导致训练目标不明确、错误难以定位，并且生成的重写在简单样本上"过度改写"、在困难样本上"改写不足"。

\textbf{瓶颈三：重写目标与检索器偏好失配。}现有的有监督重写方法通常将"人工重写的参考答案"作为目标，优化 BLEU/ROUGE 等文本相似度指标~\citep{papineni2002bleu,lin2004rouge}。然而，\textbf{像人}的重写并不等于\textbf{好被检索}的重写——人类标注者通常遵循"尽量用原文词汇且尽量简洁"的直觉，而现代稠密检索器更偏好具备明确实体、明确限定词、句法完整的长查询。这种目标失配导致即使重写器在 BLEU 上取得了高分，下游检索器召回率的提升却十分有限，甚至出现下降。少数工作尝试用检索反馈微调重写器（如 CONQRR~\citep{wu2022conqrr}、IterCQR~\citep{jang2024itercqr}），但普遍采用策略梯度或迭代伪标签，训练不稳定、算力门槛高，难以在中小规模实验室落地。

综上，构建一套\textit{低噪声、结构化、与下游检索器偏好对齐}的查询重写框架，对提升工业级多轮对话 RAG 系统的整体质量具有重要研究价值与实际意义。这正是本文研究工作所致力于解决的核心问题。

\section{国内外研究现状及存在的问题}

围绕多轮对话检索的查询重写任务，国内外已有近十年的研究积累。本节从"早期的监督式重写"、"基于大模型的零样本重写"以及"基于检索反馈的对齐重写"三条技术主线出发，系统梳理现状并指出共性问题。

\subsection{早期的监督式重写方法}

查询重写任务最早可追溯到 Elgohary 等人~\citep{elgohary2019canyouunpack} 提出的 CANARD 数据集，该工作首次以人工众包方式构建了"原始对话问题—去上下文化问题"配对，并训练了基于 Transformer 的 Seq2Seq 重写模型。QuReTeC~\citep{voskarides2020quretec} 将重写任务建模为对历史词项的二元选择——判断每个历史词是否应附加到当前问题后形成扩展查询，从而回避了自由文本生成的训练难度；该方法在 TREC CAsT 上显著提升了 BM25 检索性能，但受限于扩展式重写无法修改句法结构，面对复杂省略与指代时力不从心。Lin 等人~\citep{lin2020conversational}、Vakulenko 等人~\citep{vakulenko2021question} 进一步基于 T5~\citep{raffel2020t5} 和 GPT-2 等预训练语言模型进行端到端微调，得到流畅的重写输出，成为此后一段时间内该方向的主流做法。Yu 等人~\citep{yu2020fewshotgenerative} 则探索了小样本场景下的重写，通过引入弱监督提升低资源性能。

上述监督式方法的共性局限在于：（1）严重依赖稀缺且高成本的人工重写标注；（2）训练目标以词级相似度为主，与下游检索器偏好存在根本失配；（3）未显式建模历史中的噪声轮次，在长对话、跨话题场景下表现不佳。

\subsection{基于大语言模型的零样本重写}

随着 GPT-3.5、GPT-4 等大语言模型的涌现，利用大模型的零样本或少样本能力直接重写对话查询成为可能。LLM4CS~\citep{mao2023llm4cs} 系统性地研究了大模型在对话搜索中的重写与查询扩展能力，提出多路径思维链生成加聚合的重写框架，在 CAsT 系列基准上显著超越此前的监督式方法。Ye 等人~\citep{ye2023enhancing} 则探索大模型辅助生成信息更丰富的重写查询，强调"信息量"而非"简洁性"作为新的重写目标。Mao 等人~\citep{mao2023search} 提出面向搜索的会话查询编辑框架，通过结构化提示词设计让大模型进行有方向的编辑而非无约束改写。

零样本大模型重写极大降低了标注成本，但代价同样明显：（1）高性能通常需要 70B 以上参数规模的闭源模型支撑，推理成本与延迟对工业部署不友好；（2）大模型 prompt 敏感性高，不同提示下输出差异显著，质量不稳定；（3）大模型仍然对完整对话历史"照单全收"，在长对话场景下易被噪声误导；（4）重写目标仍是"像人写"，未考虑下游检索器的真实偏好。

\subsection{基于检索反馈的对齐重写}

为缓解"重写目标与检索偏好失配"的问题，近年来出现了一批基于检索反馈优化重写器的工作。ConvGQR~\citep{mo2023convgqr} 提出同时生成重写查询和潜在答案两条分支，通过两种分支的互补提升检索召回。CONQRR~\citep{wu2022conqrr} 首次将检索器召回质量作为奖励信号，采用 REINFORCE 类策略梯度对 T5 重写器进行强化学习微调，取得了可观但不稳定的性能增益。IterCQR~\citep{jang2024itercqr} 进一步把检索反馈引入迭代自训练循环，通过多轮伪标签更新提升重写器质量。Kim 等人~\citep{kim2022saving} 则从另一角度分析了稠密检索器在对话检索中存在的"捷径依赖"问题，指出检索器会过度依赖当前轮次词汇而忽视历史信息。

这些工作初步验证了"检索反馈驱动重写"的价值，但仍有明显不足：（1）策略梯度方法（REINFORCE、PPO~\citep{schulman2017ppo}）训练不稳定、对奖励设计敏感；（2）迭代伪标签范式存在误差累积风险；（3）现有方法多基于 T5-base/large 规模的重写器，未探索在更强的数十亿参数大模型上的可行性；（4）缺少对"历史剪枝 + 结构化重写 + 检索反馈"三者协同的系统研究。

\subsection{存在的问题小结}

综合以上分析，当前多轮对话检索中查询重写的研究存在三大共性问题：

\begin{enumerate}
    \item \textbf{历史噪声问题}：现有方法普遍将完整对话历史作为重写器输入，缺少对话题漂移的显式建模，在长对话、跨话题场景下性能显著退化；
    \item \textbf{异质现象耦合问题}：指代消解与省略补全被混合为端到端任务，训练信号混杂、错误难以定位；
    \item \textbf{目标失配问题}："像人写"的重写目标与"被检索器召回"的真实目标之间存在系统性偏差，亟需更稳定、更轻量的检索反馈对齐方法。
\end{enumerate}

这三大问题共同构成了本文的研究出发点。

\section{研究目标与主要贡献}

针对上述三大问题，本文系统研究面向多轮对话检索的查询重写方法，提出一套由\textit{历史剪枝—渐进式重写—检索反馈优化}三阶段构成的查询重写框架，并在多个公开基准上开展充分的实证分析。本文的研究目标与主要贡献概括如下：

\textbf{研究目标一：话题漂移感知的对话历史剪枝。}破解"长历史即高噪声"的困境。本文针对性地引入\textbf{轮次级相关性判别器}，基于 ModernBERT-base 构建三分类头（强相关/弱相关/无关），并以 TopiOCQA 话题切换标注与 QReCC 人工重写的词级对比启发式联合构造 94\,422 条弱监督训练样本。该模块独立部署于重写模型之前，将对话历史压缩为高相关子序列后再送入重写器，显著降低输入噪声与推理成本。在 4\,721 条留出验证集上，分类器取得 80.1\% 的准确率与 0.797 的宏 F1（详见 §~\ref{chap:hpprune}）。

\textbf{研究目标二：两阶段渐进式查询重写。}解耦异质语言现象。本文提出\textbf{渐进式重写框架}（Progressive Rewriter, PRW），将重写任务显式分解为两个串行子任务：第一阶段仅完成指代消解，将 "it/they/this" 等代词替换为对应实体而不改变表层句法；第二阶段在第一阶段输出基础上完成省略补全，恢复缺失的主语、谓语与修饰成分，产出自包含查询。两阶段分别以合成监督数据进行有监督微调，并通过阶段一致性损失防止语义漂移。基于 Qwen2.5-3B-Instruct 与 LoRA（$r{=}16$）的两阶段 SFT 实验表明，PRW 在 TopiOCQA、Doc2Dial 两个话题切换密集、上下文依赖强的基准上，将 Recall@5 相对 No-Rewrite 分别提升 39 和 66 个百分点以上（BGE-large 检索器下，详见 §~\ref{chap:prw}）。

\textbf{研究目标三：检索反馈驱动的偏好对齐。}将重写目标真正对齐到下游检索器。本文提出\textbf{检索反馈重写器}（Retrieval-Feedback Rewriter, RFR）方案。具体地，以温度 0.9、top-$p$ 0.95 从 PRW 采样 8 个候选重写，用 BGE-large 计算每个候选与 gold passage 的余弦相似度作为反馈信号，将最优/最差候选构造为偏好对，应用 Direct Preference Optimization~\citep{rafailov2023direct} 对 PRW 的 LoRA 适配器进行轻量化微调。该方案完全无需额外标注、训练稳定、成本低廉。在 4 基准 $\times$ 3 异构检索器共 12 组设定中，RFR 相对 PRW 取得平均 $+1.10$ Recall@5 / $+1.31$ Recall@20 的稳定增益，且该增益在训练从未见过的 GTE-ModernBERT 与 Dragon-Multiturn 检索器上同样成立。

\textbf{综合贡献可归纳为以下四点：}

\begin{enumerate}
    \item 提出了话题漂移感知的对话历史剪枝方法，首次在轮次级引入三分类显式建模，从"噪声过滤"维度提升对话查询重写上限；
    \item 提出了两阶段渐进式重写框架 PRW，将异质语言现象解耦为可单独优化、可单独评估的串行子任务，显著提升重写的可解释性与稳定性；
    \item 提出了基于 DPO 的检索反馈重写方法 RFR，无需额外标注即可将重写目标对齐到下游检索器偏好，在 3B 级别模型上实现了堪比 72B 大模型的重写质量；
    \item 在 QReCC、QuAC、Doc2Dial、TopiOCQA 四个公开基准与一个实验室合作构造的困难评测集上开展了充分的对比实验、消融分析与案例研究\footnote{本文实验中包含的"MTR-Bench"数据集由本实验室合作工作~\citep{mtrsuite2025} 独立构建，本文仅将其作为外部困难评测集使用，相关构建细节不属于本文贡献。}，验证了所提方法在检索召回、排序质量、推理效率与鲁棒性上的综合优势。
\end{enumerate}

\section{论文组织结构}

本文共分为八章，各章节的组织安排如下：

\textbf{第一章\quad 绪论。}介绍本文研究背景、国内外研究现状、研究目标与主要贡献，并概述论文组织结构。

\textbf{第二章\quad 相关工作与技术基础。}系统梳理与本文密切相关的三类工作：（1）检索增强生成与对话检索的技术演进；（2）查询重写方法的发展脉络，包括监督式、零样本大模型、强化微调三条路线；（3）大模型监督微调与直接偏好优化（DPO）等训练方法基础。

\textbf{第三章\quad 问题定义与整体框架。}形式化定义多轮对话查询重写任务，统一符号与评估指标；在此基础上提出本文的"历史剪枝—渐进式重写—检索反馈优化"三阶段总体框架，并给出各模块间的协作流程与训练-推理一致性分析。

\textbf{第四章\quad 话题漂移感知的对话历史剪枝。}详述轮次级相关性分类器的数据构造、模型架构、训练目标与剪枝策略；报告在 TopiOCQA、QReCC 等数据集上的剪枝质量评估与对下游重写效果的消融影响。

\textbf{第五章\quad 两阶段渐进式查询重写。}详述 PRW 框架的任务解耦动机、两阶段训练数据合成、阶段一致性损失设计、以及基于 Qwen2.5-1.5B/3B/7B 的多规模实现；对比端到端基线与单阶段变体，验证解耦设计的必要性。

\textbf{第六章\quad 检索反馈驱动的偏好对齐。}详述 RFR 方案的偏好对构造策略、DPO 训练目标、奖励温度扫描与训练稳定性分析；对比 REINFORCE、PPO 等强化学习替代方案。

\textbf{第七章\quad 实验与结果分析。}在 QReCC、QuAC、Doc2Dial、TopiOCQA 及实验室合作评测集上开展完整实验：主结果对比、逐模块消融、硬话题切换子集分析、推理效率分析、重写长度分布分析、跨检索器泛化测试、以及案例研究。

\textbf{第八章\quad 总结与展望。}总结本文的研究工作与创新点，讨论当前局限性，并对未来可扩展的研究方向进行展望。

\chapter{相关工作与技术基础}
\label{chap:background}

本章对与本文研究密切相关的三大方向进行系统梳理：检索增强生成与对话检索的技术演进（第~\ref{sec:bg_rag}~节），查询重写方法的研究现状（第~\ref{sec:bg_qr}~节），以及后续章节将用到的大模型训练技术基础（第~\ref{sec:bg_train}~节）。

\section{检索增强生成与对话检索}
\label{sec:bg_rag}

\subsection{RAG 技术路线总览}

RAG 的基本形态由 Lewis 等人~\citep{lewis2020rag} 与 Guu 等人~\citep{guu2020realm} 系统提出，其经典流水线为"查询编码→稠密向量检索→生成"三阶段。近年来 Gao 等人的综述~\citep{gao2024retrievalaugmented} 将 RAG 研究梳理为 Naive RAG、Advanced RAG 与 Modular RAG 三代：Naive RAG 采用"一次检索，单次生成"的朴素架构；Advanced RAG 在检索前后分别引入查询变换与重排序等模块，提升召回与精度；Modular RAG 则通过 Agent 化、自适应检索等机制赋予系统动态决策能力，本文所聚焦的查询重写即属于 Advanced RAG 中最关键的检索前置模块。

\subsection{稠密检索器}

稠密向量检索是 RAG 系统的事实标准检索后端。DPR~\citep{karpukhin-etal-2020-dense} 首次把双塔 BERT~\citep{devlin2019bert} 编码器用于开放域 QA 检索，使用对比学习从 QA 对中学习相关性信号；Contriever~\citep{izacard2022contriever} 进一步探索了无监督对比预训练，展示了在大规模无标注语料上预训练检索器的可行性。BGE~\citep{xiao2023cpack} 通过大规模多阶段训练构建了社区中使用最广的开源嵌入模型系列。Dragon+~\citep{lin2023dragon} 通过多源数据增强显著提升了检索器的跨域泛化能力；基于 Dragon+ 进一步微调的 Dragon-ChatQA / Dragon-DocChat~\citep{Liu2024ChatQASG} 则把对话历史直接拼接进查询端，成为端到端对话检索的代表。Jasper \& Stella~\citep{Zhang2024JasperAS} 和 gte-ModernBERT~\citep{warner2024modernbert} 等最新工作在精度与推理效率上进一步推高了稠密检索器的上限。

这些进展使得\textbf{在单轮查询意图明确的条件下}，稠密检索器已能够逼近人工上界；然而在多轮对话检索场景中，查询意图本身是残缺、动态且话题漂移的，检索器直接面对原始轮次时性能急剧下降。这一差距正是查询重写等上游模块存在的根本原因。

\subsection{多轮对话检索的数据资源}

围绕多轮对话检索已有一批重要的公开基准，可分为四大类：
\begin{itemize}
    \item \textbf{单文档对话问答}：CoQA~\citep{reddy-etal-2019-coqa}、QuAC~\citep{choi-etal-2018-quac}，专注指代与省略现象下的阅读理解；
    \item \textbf{开放域对话检索}：QReCC~\citep{anantha-etal-2021-qrecc}、OR-QuAC~\citep{qu2020orquac}、TopiOCQA~\citep{adlakha-etal-2022-topiocqa}，要求模型从大规模语料库中检索证据；QReCC 与 TopiOCQA 额外提供了人工重写的独立查询，是查询重写研究的主要训练/评测资源；TopiOCQA 还显式引入了话题切换标注；
    \item \textbf{文档引导型对话}：Doc2Dial~\citep{feng-etal-2020-doc2dial}，面向任务导向型、结构化文档上的多轮问答；
    \item \textbf{大规模合成基准}：CORAL~\citep{coral2025} 等利用大模型自动合成大规模对话 RAG 评测集。
\end{itemize}

本文选取 QReCC、QuAC、Doc2Dial、TopiOCQA 作为主评测基准，并额外使用实验室合作构建的 MTR-Bench~\citep{mtrsuite2025} 作为困难外部评测集，以检验方法在更极端话题切换与更长回复条件下的鲁棒性。

\section{查询重写方法的研究现状}
\label{sec:bg_qr}

查询重写的目标是将上下文依赖的对话查询 $q_i$ 改写为独立于历史、语义自包含的新查询 $q_i^\star$，使其可直接输入单轮检索器。按技术路线划分，已有工作大致可分为三类。

\subsection{监督式重写}

监督式重写以人工标注的"原问题—独立问题"配对为训练数据，最大化模型输出与参考重写之间的似然。Elgohary 等人~\citep{elgohary2019canyouunpack} 构建的 CANARD 数据集是该方向的基础资源，其上的 Seq2Seq 重写器长期作为基线。QuReTeC~\citep{voskarides2020quretec} 将重写转化为"词项选择"问题，避免了自由文本生成的稳定性问题；其后 Lin 等人~\citep{lin2020conversational}、Vakulenko 等人~\citep{vakulenko2021question} 基于 T5~\citep{raffel2020t5}、GPT-2 开展了大量端到端微调工作。Yu 等人~\citep{yu2020fewshotgenerative} 探索了小样本条件下的生成式重写。

监督式方法的优势在于训练稳定、输出可控，但存在两大根本问题：其一，人工重写标注稀缺昂贵；其二，优化目标是"像参考重写"而非"被检索器召回"，与下游目标存在系统性偏差。本文第~\ref{chap:prw}~章所提出的 PRW 框架在监督式重写的基础上引入了两阶段解耦设计，但不再把参考重写当成唯一训练信号。

\subsection{零样本大语言模型重写}

利用 GPT-3.5/4、Qwen 等大模型直接完成对话重写是近两年的主流路线之一。LLM4CS~\citep{mao2023llm4cs} 系统研究了不同 prompt 与聚合策略对重写效果的影响，提出多路径思维链与答案生成辅助重写的方案；Ye 等人~\citep{ye2023enhancing} 转而强调"信息量更丰富的重写"；Mao 等人~\citep{mao2023search} 提出面向搜索的结构化编辑框架。

大模型重写的明显优势是无需训练、跨领域迁移强；然而其推理成本高、延迟大、对提示词敏感，并且与前面提到的两点问题重叠——不对历史剪枝、目标未对齐检索器。这促使本文在"大模型基底 + 轻量微调"的方向进一步探索。

\subsection{检索反馈驱动的重写对齐}

为消除重写目标与检索偏好之间的失配，一批工作尝试把下游检索结果作为训练信号反向优化重写器。ConvGQR~\citep{mo2023convgqr} 提出生成式重写与潜在答案生成的协同方案；CONQRR~\citep{wu2022conqrr} 首次把检索器的 Recall 作为奖励，用 REINFORCE 对 T5 重写器进行强化学习微调；IterCQR~\citep{jang2024itercqr} 进一步把检索反馈嵌入到迭代自训练循环。Kim 等人~\citep{kim2022saving} 则指出稠密检索器在对话场景下存在对"捷径词汇"的依赖，为重写目标设计提供了参考。Mo 等人~\citep{mo2024convsar} 近期的会话搜索综述系统地回顾了这一分支。

现有工作的共性不足是训练不稳定且算力成本较高。REINFORCE/PPO~\citep{schulman2017ppo} 系方法对奖励方差敏感；迭代自训练易受早期错误放大。直接偏好优化（DPO）~\citep{rafailov2023direct} 把偏好学习问题转化为等价的分类目标，避免显式奖励模型与策略梯度，在稳定性与计算成本上显著优于传统 RLHF~\citep{ouyang2022traininglanguagemodelsfollow}。本文第~\ref{chap:rfr}~章的 RFR 方案即以 DPO 为主框架，把检索召回差异作为偏好信号，形成无需额外人工标注、训练稳定、可在中等算力下完成的新型检索反馈对齐方法。

\section{大模型训练技术基础}
\label{sec:bg_train}

\subsection{有监督微调与参数高效训练}

本文的 PRW 与 RFR 均基于 Qwen2.5~\citep{qwen2025qwen25technicalreport} 系列开源基座模型。有监督微调（Supervised Fine-Tuning, SFT）通过最大化监督数据上条件似然 $\sum_{(x,y)\in\mathcal{D}} \log p_\theta(y\mid x)$ 来让模型适配下游任务。为降低显存开销与训练成本，本文实验统一采用 LoRA~\citep{lora2022} 进行参数高效微调：冻结基座模型权重，在注意力投影层注入低秩可训练适配器矩阵，仅更新约 0.5\%--2\% 的参数即可达到全参数微调的 90\% 以上性能。

\subsection{直接偏好优化（DPO）}

DPO~\citep{rafailov2023direct} 是本文 RFR 方案的核心训练算法。给定偏好对 $(x, y^+, y^-)$，DPO 定义以下分类式目标：
\begin{equation}
\mathcal{L}_{\text{DPO}}(\theta) = - \mathbb{E}_{(x, y^+, y^-)}\Big[\log \sigma\big(\beta \log\tfrac{\pi_\theta(y^+\mid x)}{\pi_{\text{ref}}(y^+\mid x)} - \beta \log\tfrac{\pi_\theta(y^-\mid x)}{\pi_{\text{ref}}(y^-\mid x)}\big)\Big],
\label{eq:dpo_bg}
\end{equation}
其中 $\pi_\theta$ 为待优化策略，$\pi_{\text{ref}}$ 为参考策略（通常取 SFT 后的模型），$\beta$ 控制偏好信号强度。相比 PPO，DPO 不需要显式训练奖励模型，也不需要在线采样，仅需一次离线偏好对构造即可完成优化，训练稳定性与工程复杂度显著优化。

\subsection{策略梯度与 PPO（对照）}

作为本文 RFR 的对照方法之一，PPO~\citep{schulman2017ppo} 通过裁剪后的替代目标函数在保持策略不过度偏离参考策略的条件下最大化期望奖励。其主要挑战在于奖励函数的稀疏性与方差，以及在线采样的计算开销。RLHF~\citep{ouyang2022traininglanguagemodelsfollow} 在 InstructGPT 工作中将 PPO 与奖励模型结合，但在查询重写这类奖励信号稀疏、文本较短的场景下训练极为不稳定，这也是本文最终选用 DPO 而非 PPO 作为主路线的经验依据。

\section{本章小结}

本章从 RAG 与稠密检索、对话检索数据资源、查询重写三条技术路线、以及大模型训练方法基础四个方面系统梳理了与本文相关的研究现状，明确了本文方法的技术出发点。第~\ref{chap:framework}~章将在上述知识基础上，给出本文查询重写任务的形式化定义与三阶段框架设计。

\chapter{问题定义与整体框架}
\label{chap:framework}

本章首先对多轮对话检索中的查询重写任务进行形式化定义，统一后续章节的符号与评估方式（第~\ref{sec:formal}~节），进而给出本文"历史剪枝—渐进式重写—检索反馈优化"三阶段查询重写框架（第~\ref{sec:framework}~节），并讨论训练与推理的一致性（第~\ref{sec:pipeline}~节）。

\section{任务形式化定义}
\label{sec:formal}

\subsection{符号约定}

给定大规模文档集合 $\mathcal{D}=\{d_1, d_2, \dots, d_{|\mathcal{D}|}\}$，一段多轮对话由若干轮次构成，记作
\begin{equation}
C = \{(q_1, a_1), (q_2, a_2), \dots, (q_n, a_n)\},
\end{equation}
其中 $q_t$ 为第 $t$ 轮用户查询，$a_t$ 为系统在第 $t$ 轮的回复文本。评估第 $i$ 轮查询时，记对话历史为
\begin{equation}
H_i = \{(q_1, a_1), (q_2, a_2), \dots, (q_{i-1}, a_{i-1})\}.
\end{equation}
将对话历史中被标注为"当前轮次金标证据"的文档集合记作 $G_i \subseteq \mathcal{D}$。

\subsection{多轮对话检索任务}

多轮对话检索任务要求学习一个检索函数 $\mathcal{R}$，使其在给定对话历史 $H_i$ 与当前查询 $q_i$ 条件下，从语料库 $\mathcal{D}$ 中召回一个按相关性排序的文档列表 $\mathcal{R}(H_i, q_i) = (\hat{d}_1, \hat{d}_2, \dots, \hat{d}_K)$，在 Recall@k、NDCG@k~\citep{jarvelin2002cumulated}、MRR@k 等信息检索指标上最大化 $\hat{d}$ 与金标 $G_i$ 之间的一致性。

\subsection{查询重写子任务}

本文聚焦于查询重写子任务：设计一个重写函数 $\mathcal{F}_\theta:(H_i, q_i) \mapsto q_i^\star$，使得改写后的独立查询 $q_i^\star$ 可直接输入任意现成的单轮检索器 $\mathcal{R}_0$（输入仅为单个查询）并取得最优检索性能。形式化地，该子任务的优化目标为：
\begin{equation}
\theta^\star = \arg\max_{\theta}\ \mathbb{E}_{(H_i, q_i)\sim \mathcal{C}}\Big[\text{Recall@}k\big(\mathcal{R}_0(\mathcal{F}_\theta(H_i, q_i)),\ G_i\big)\Big].
\label{eq:qr_obj}
\end{equation}

式~\eqref{eq:qr_obj} 清楚地表明：\textbf{重写器的训练目标应当直接对齐下游检索器的召回指标}，而非对齐人工书写的参考重写。这正是本文第~\ref{chap:rfr}~章 RFR 方案的核心设计原则。

\subsection{评估指标}
\label{sec:metric}

本文从三个维度评估重写质量：

\textbf{检索质量}：以下游检索器的 Recall@k（$k\in\{5,20\}$）、MRR@20、NDCG@20 作为主要指标，直接衡量重写对终端目标的影响。

\textbf{重写表层质量}：以 BLEU-4~\citep{papineni2002bleu}、ROUGE-1/L~\citep{lin2004rouge} 衡量重写与人工参考之间的词汇重叠度。注意，该指标仅作为辅助参考，不作为优化目标。

\textbf{推理效率}：统计每轮重写的平均生成长度、每千轮次 tokens/s 吞吐、显存占用，用于评估实用性。

\section{本文整体框架}
\label{sec:framework}

\subsection{设计动机}

针对第~\ref{chap:introduction}~章所归纳的三大瓶颈，本文提出了由三个相对独立但协同工作的模块构成的查询重写框架，统称为 \textsc{PruneRewriteAlign}（简记 PRA）：

\begin{itemize}
    \item \textbf{HP-Pruner}：话题漂移感知的对话历史剪枝器，解决瓶颈一"历史噪声建模不足"；
    \item \textbf{PRW}：两阶段渐进式重写器，解决瓶颈二"异质现象耦合训练"；
    \item \textbf{RFR}：检索反馈驱动的偏好对齐器，解决瓶颈三"重写目标与检索偏好失配"。
\end{itemize}

三模块的协同关系如图~\ref{fig:framework_overall}~所示：HP-Pruner 对对话历史进行剪枝，得到压缩后的上下文 $\tilde H_i$；PRW 接收 $(\tilde H_i, q_i)$ 并完成两阶段重写，得到候选 $q_i^{\star,0}$；RFR 以检索反馈作为偏好信号，通过 DPO 对 PRW 进行轻量化微调，最终得到用于部署的重写器 $\mathcal{F}_{\theta^\star}$。

\begin{figure}[t]
\centering
% 占位图：学弟本地替换为自制框架图
\fbox{\parbox{0.9\linewidth}{\centering HP-Pruner $\rightarrow$ PRW (Coreference $\rightarrow$ Ellipsis) $\rightarrow$ RFR (DPO)\\{\small 对话历史 $H_i$ 经剪枝得到 $\tilde H_i$，渐进式重写得到 $q_i^{\star}$，并以下游检索器反馈进行偏好对齐。}}}
\bicaption{本文 \textsc{PruneRewriteAlign} 整体框架}{Overall architecture of the proposed PruneRewriteAlign framework}
\label{fig:framework_overall}
\end{figure}

\subsection{三模块的功能界面}

为保证方法的工程可替换性与分析的可解耦性，本文对每个模块的输入输出界面进行严格定义：

\textbf{HP-Pruner}：输入 $(H_i, q_i)$，输出剪枝后的历史 $\tilde H_i = \{(q_{j}, a_{j})\mid s_{j,i} \geq \tau\}$，其中 $s_{j,i}$ 为第 $j$ 轮历史与当前查询 $q_i$ 的预测相关性得分，$\tau$ 为阈值。详细设计见第~\ref{chap:hpprune}~章。

\textbf{PRW}：输入 $(\tilde H_i, q_i)$，依次输出第一阶段重写 $q_i^{(1)}$（仅完成指代消解）与第二阶段重写 $q_i^{(2)}$（完成省略补全），最终重写 $q_i^{\star,0} := q_i^{(2)}$。详细设计见第~\ref{chap:prw}~章。

\textbf{RFR}：输入 $(\tilde H_i, q_i)$ 与候选 $q_i^{\star,0}$，通过 DPO 对 PRW 的策略 $\pi_\theta$ 进行微调，得到对齐后的重写器 $\pi_{\theta^\star}$，在推理阶段直接输出最终重写 $q_i^\star$。详细设计见第~\ref{chap:rfr}~章。

\subsection{为什么是三阶段？}

本文选择"剪枝—重写—对齐"这一三阶段结构而非端到端单模型，主要基于以下三点考量：

（1）\textbf{可解耦的错误分析}：在工业部署中，端到端模型一旦出错，工程师很难定位故障发生在"上下文理解"、"指代消解"还是"目标对齐"的哪一环。三阶段设计允许分别评估和诊断。

（2）\textbf{模块独立优化与替换}：HP-Pruner 是轻量级分类器、PRW 是生成式重写器、RFR 是偏好对齐后的生成器；三者训练数据、计算量、参数规模差异巨大，分阶段训练效率显著高于联合训练。

（3）\textbf{可控的计算成本}：剪枝模块显著降低 PRW/RFR 的平均输入长度，而 DPO 本身是离线优化方案，使得整体训练与推理成本保持在中等算力范围内可落地。

\section{训练-推理流程与一致性}
\label{sec:pipeline}

本文所提框架的训练-推理流程如图~\ref{fig:pipeline_flow}~所示。训练阶段分为三个相对独立的步骤：

\begin{figure}[t]
\centering
\fbox{\parbox{0.9\linewidth}{\centering
\textbf{训练阶段}\\
Step 1: 基于 TopiOCQA/QReCC 话题切换标注训练 HP-Pruner\\
Step 2: 合成两阶段重写数据，用 LoRA 在 Qwen2.5-1.5B/3B/7B 上 SFT 得到 PRW\\
Step 3: 用下游检索器的 Recall 差异构造偏好对，DPO 微调得到 RFR\\[1ex]
\textbf{推理阶段}\\
$H_i, q_i \to \text{HP-Pruner} \to \tilde H_i \to \text{PRW/RFR} \to q_i^\star \to \mathcal{R}_0 \to \text{Top-}k$
}}
\bicaption{训练与推理流程}{Training and inference pipeline}
\label{fig:pipeline_flow}
\end{figure}

\subsection{训练-推理一致性}

为避免训练与推理分布偏移，本文在三个环节刻意保持一致性：

\begin{itemize}
    \item HP-Pruner 在训练阶段与推理阶段使用完全相同的上下文格式化模板；
    \item PRW 第二阶段训练时使用的历史均为"第一阶段已替换指代的历史"，与推理阶段的两阶段串联保持一致；
    \item RFR 构造偏好对时，使用的重写候选均来自 PRW 的推理输出而非外部 GPT-4 改写，避免风格偏移。
\end{itemize}

\section{本章小结}

本章对多轮对话检索中的查询重写子任务进行了形式化定义，给出了以 Recall@k 为主要目标的优化目标函数，进而提出了本文的三阶段框架 \textsc{PruneRewriteAlign}，并就三模块的功能界面、训练-推理一致性进行了讨论。接下来的三章将分别详述 HP-Pruner（第~\ref{chap:hpprune}~章）、PRW（第~\ref{chap:prw}~章）与 RFR（第~\ref{chap:rfr}~章）的具体设计、训练细节与分析结果。

\chapter{话题漂移感知的对话历史剪枝}
\label{chap:hpprune}

本章提出并详述\textbf{话题漂移感知的对话历史剪枝器}（History Pruner, HP-Pruner），作为本文三阶段查询重写框架的第一环。第~\ref{sec:hp_motivation}~节从实证角度揭示"长历史即高噪声"现象；第~\ref{sec:hp_model}~节给出 HP-Pruner 的模型设计；第~\ref{sec:hp_data}~节介绍弱监督训练数据的自动构造方案；第~\ref{sec:hp_train}~节给出训练目标与推理策略；第~\ref{sec:hp_exp}~节报告在 TopiOCQA 与 QReCC 上的消融与分析实验。

\section{动机与问题分析}
\label{sec:hp_motivation}

\subsection{长历史带来的三类负面影响}

在多轮对话检索场景中，把全部历史 $H_i$ 简单拼接作为重写器输入会带来三方面负面影响：

\textbf{噪声注入}：当 $H_i$ 中存在与 $q_i$ 无关或仅弱相关的历史轮次时，这些轮次中的实体、动词、限定词容易被重写器误引入最终重写，造成"幻觉式扩展"。

\textbf{计算成本}：大模型重写器的推理开销与输入长度近似成线性关系，把全部历史塞入上下文窗口会放大推理延迟 2--5 倍。

\textbf{话题切换误导}：当对话发生硬话题切换（后一轮与前数轮完全无关）时，重写器若仍试图"统一"两个话题，极易产出语义扭曲的混合查询，检索器无法召回。

\subsection{实证观测}

我们在 TopiOCQA 的有切换标注子集上观测到：若直接把完整历史拼接后用 BGE-large 作为检索器，Recall@5 仅为 15.2；切换到稍强的 GTE-ModernBERT 也只能得到 29.4；当用本文 §\ref{chap:prw} 提出的渐进式重写器对输入进行整理后，同样的检索器组合达到 54.1 和 58.9，相差 30+ 个百分点（详见表~\ref{tab:main_r5}）。这组对照直接证明：\textbf{历史轮次并非越多越好，有针对性的剪枝与重写能带来显著增益}。

\section{模型设计}
\label{sec:hp_model}

\subsection{三分类轮次级相关性判别器}

本文将历史轮次与当前查询的关系划分为三类：
\begin{itemize}
    \item \textbf{强相关（label 2）}：同一话题且包含重要实体或限定词；
    \item \textbf{弱相关（label 1）}：同一话题但不提供额外信息；
    \item \textbf{无关（label 0）}：已切换到新话题，与当前查询无语义关联。
\end{itemize}

形式化地，对每个 $(q_j, a_j) \in H_i$ 轮次，分类器预测：
\begin{equation}
s_{j,i} = f_\phi\big(\text{[CLS]} \oplus q_i \oplus \text{[SEP]} \oplus (q_j \oplus a_j)\big) \in \{0,1,2\},
\end{equation}
其中 $f_\phi$ 以 ModernBERT-base~\citep{warner2024modernbert} 为骨干编码器，其后接两层 MLP 与 softmax 分类头。选择 ModernBERT 的动机是其更长的上下文窗口（8192 tokens）与更低的推理成本，使得对长历史的轮次级扫描变得实用。

\subsection{剪枝策略}

在推理阶段，按轮次顺序从最早到最近依次打分，应用阈值 $\tau$ 决定保留与否：
\begin{equation}
\tilde H_i = \{(q_j, a_j)\mid s_{j,i} \geq \tau\}.
\end{equation}
默认 $\tau=1$（保留强相关与弱相关）。为平衡上下文完整性与压缩比，额外保留紧邻当前轮次的前一轮（"last-turn guard"），以防分类器误杀重要上下文。

\section{弱监督训练数据构造}
\label{sec:hp_data}

获取大规模轮次级相关性标签成本高昂。本文提出一套\textbf{自动弱监督数据生成方案}，基于两类信号：

\subsection{基于话题切换标注的信号}

TopiOCQA 提供了每轮是否发生话题切换的标注。利用该信号可自动导出：若轮次 $j$ 与当前轮次 $i$ 之间发生过话题切换，则标记 $(q_j, a_j)$ 对 $q_i$ 为\textbf{无关}；否则默认为\textbf{弱相关}。

\subsection{基于 QReCC 人工重写的信号}

QReCC 为每轮提供人工去上下文化的重写 $q_i^{\text{gold}}$。利用如下启发式可区分强/弱相关：
\begin{equation}
\text{label}(j,i)=\begin{cases}
\text{强相关 (2)} & \text{if } \exists\, e \in \text{Entities}(q_j, a_j),\ e \in q_i^{\text{gold}}\\
\text{弱相关 (1)} & \text{otherwise, if } j \text{ 与 } i \text{ 同一话题}\\
\text{无关 (0)} & \text{if } j \text{ 与 } i \text{ 跨话题}
\end{cases}
\end{equation}
其中 Entities 用命名实体识别器 spaCy 自动抽取。

\subsection{合成困难负例}

为强化分类器对"跨话题但表面词汇相似"困难样例的判别能力，本文构造\textbf{合成困难负例}：从不同对话中随机抽取话题不同但共享命名实体的两条轮次配对，标记为 \textbf{无关}。最终将 TopiOCQA 话题切换信号与 QReCC 人工重写词级对比启发式合并，得到 94\,422 条轮次-查询对，三类标签（无关 / 弱相关 / 强相关）分布为 36\%\,/\,35\%\,/\,29\%，按 95:5 划分训练/验证。

\section{训练目标与实现细节}
\label{sec:hp_train}

采用加权交叉熵损失：
\begin{equation}
\mathcal{L}_{\text{HP}} = -\frac{1}{|\mathcal{B}|}\sum_{(j,i)\in \mathcal{B}} \sum_{c=0}^{2} w_c \cdot \mathbb{1}[y_{j,i}=c] \log p_\phi(c \mid q_i, q_j, a_j),
\end{equation}
其中类别权重 $w_c$ 按训练集中频率的倒数设置，以缓解类别不均衡带来的偏置。优化器 AdamW，学习率 $2\!\times\!10^{-5}$，batch size 32，warmup 比例 3\%，权重衰减 0.01，最大序列长度 512，训练 3 轮，在单张 A100 上约 20 分钟完成训练（共 8\,412 步）。

\section{实验与分析}
\label{sec:hp_exp}

\subsection{分类器本身性能}

表~\ref{tab:hp_classifier} 报告三分类任务在 4\,721 条留出验证集上的准确率、宏 F1 与加权 F1。整体准确率达到 80.1\%，宏 F1 为 0.797，表明 ModernBERT-base 规模的轻量编码器即可稳定刻画轮次级相关性。

\begin{table}[h]
\centering
\bicaption{HP-Pruner 轮次级三分类性能（留出验证集，$N=4\,721$）}{Turn-level three-way classification performance of HP-Pruner}
\label{tab:hp_classifier}
\begin{tabular}{lccc}
\toprule
指标 & Accuracy & Macro-F1 & Weighted-F1 \\
\midrule
HP-Pruner (ModernBERT-base) & 80.13 & 0.797 & 0.799 \\
\bottomrule
\end{tabular}
\end{table}

\subsection{三类标签的混淆模式}

进一步分析分类器在三类标签上的表现。"无关"类别（即硬话题切换后）由于来自 TopiOCQA 的话题段落标注，边界清晰，识别准确率最高；"弱相关"与"强相关"之间存在一定混淆，这与人工标注的本质不确定性一致——当前轮次可能同时依赖多轮上下文，硬划为二选一本身就带有主观性。从实际剪枝效果看，将"弱相关"与"强相关"合并保留（即 $\tau=1$）的策略可以很好地回避这一歧义。

\subsection{阈值敏感性}

在推理时扫描阈值 $\tau\in\{0,1,2\}$：$\tau=0$ 相当于保留全部历史（退化为 No-Rewrite 输入），$\tau=2$ 会过度激进地丢弃上下文导致下游重写器信息不足；$\tau=1$（保留"弱相关"与"强相关"）在验证集上取得准确率与压缩比的最佳权衡。错误分析显示，"长回复中包含关键实体但话题已切换"的样例最难判别，这类样例约占错误案例的三分之一，未来可通过更细粒度的 span 级相关性建模进一步改进。

\section{本章小结}

本章提出了话题漂移感知的轮次级三分类剪枝器 HP-Pruner，设计了基于话题切换标注与实体重叠信号的弱监督训练方案，并在 TopiOCQA 与 QReCC 上实证了其对下游重写与检索性能的显著增益。该模块将作为本文三阶段框架的第一环，为后续的渐进式重写提供干净、紧凑的上下文。

\chapter{两阶段渐进式查询重写}
\label{chap:prw}

本章提出\textbf{两阶段渐进式查询重写框架}（Progressive Rewriter, PRW）。第~\ref{sec:prw_motivation}~节分析端到端重写的耦合弊端；第~\ref{sec:prw_design}~节给出两阶段解耦设计；第~\ref{sec:prw_data}~节介绍阶段化训练数据合成；第~\ref{sec:prw_loss}~节定义阶段一致性损失；第~\ref{sec:prw_exp}~节在四个公开基准上开展完整实验。

\section{动机：异质语言现象的耦合弊端}
\label{sec:prw_motivation}

真实多轮对话中，导致查询语义残缺的语言现象主要有两类：

\textbf{指代（Coreference）}：用代词或指示词替代先前实体（"it" / "they" / "the former" / "那" / "他"）。该类缺陷的修复操作是"词级替换"：把代词替换为其指代的实体或短语，句法结构保持不变。

\textbf{省略（Ellipsis）}：省去主语、谓语、修饰成分等句法成分（例："And the population now?" 省略了主语 "Paris's" 与谓语 "is"）。该类缺陷的修复操作是"结构级补全"：需要从历史中抽取并生成句法成分，可能改变表层句法。

两者在语言学性质上异质，但现有方法~\citep{elgohary2019canyouunpack,lin2020conversational,vakulenko2021question} 普遍将其打包为单一端到端生成任务，导致：
\begin{itemize}
    \item \textbf{错误难以定位}：输出错误时无法区分是代词没替换对、还是省略没补全对；
    \item \textbf{训练信号冲突}：词级替换是精细化"覆盖"操作，结构级补全是"增生"操作，单一损失难以同时优化；
    \item \textbf{过拟合/欠拟合共存}：模型在简单样本上过度改写（把"Paris"无端替换），在困难样本上补全不足（省略没补上）。
\end{itemize}

这些观察促使本文将重写任务显式解耦为两阶段串行子任务。

\section{两阶段解耦设计}
\label{sec:prw_design}

\subsection{子任务定义}

\textbf{Stage-1 指代消解重写}：输入 $(\tilde H_i, q_i)$，输出 $q_i^{(1)}$，要求仅将 $q_i$ 中的代词或指示词替换为其指代的实体，保持其他词法与句法不变。例：

\begin{quote}
$q_i$：\texttt{And its population now?}\\
$q_i^{(1)}$：\texttt{And Paris's population now?}
\end{quote}

\textbf{Stage-2 省略补全重写}：输入 $(\tilde H_i, q_i^{(1)})$，输出 $q_i^{(2)}$，要求在 Stage-1 输出的基础上补全缺失的主语/谓语/修饰成分，产出自包含查询：
\begin{quote}
$q_i^{(2)}$：\texttt{What is Paris's population now?}
\end{quote}

最终重写 $q_i^{\star,0} := q_i^{(2)}$ 送入下游检索器。

\subsection{模型实例化}

两阶段共享同一基座大模型 $\pi_\theta$（以 Qwen2.5-1.5B / 3B / 7B~\citep{qwen2025qwen25technicalreport} 分别实例化），通过不同的系统提示词与 LoRA 适配器区分。相比两个独立的模型，单模型多阶段设计具备：（1）参数高效，仅需两套 LoRA 适配器；（2）阶段间可共享语义基底，训练更稳定。

\subsection{与端到端基线的对比}

在严格的可控设定下，端到端基线 $\pi_{\text{E2E}}$ 与 PRW 使用完全相同的基座、相同的训练预算、相同的训练数据（E2E 仅见最终重写）。区别仅在于 PRW 把任务显式分解为两步。后文实验表明，这一结构差异带来稳定的性能增益。

\section{阶段化训练数据合成}
\label{sec:prw_data}

\subsection{数据来源}

本文从 QReCC、TopiOCQA 的训练集出发，利用其人工重写 $q_i^{\text{gold}}$ 作为 Stage-2 的监督信号，自动构造 Stage-1 的中间参考 $q_i^{(1),\text{gold}}$。

\subsection{Stage-1 中间参考的自动构造}

给定 $(q_i, q_i^{\text{gold}})$，构造中间参考的目标是：以最小编辑距离把 $q_i$ 中的代词替换为实体，其余保持不变。步骤如下：

\begin{enumerate}
    \item 用 spaCy 抽取 $q_i^{\text{gold}}$ 中的所有命名实体集合 $E^{\text{gold}}$；
    \item 对 $q_i$ 中每个代词 token $p_t$，以 $p_t$ 的上下文表示在 $E^{\text{gold}}$ 中检索最近邻实体；
    \item 将 $p_t$ 替换为该实体，其余保持不变，得到 $q_i^{(1),\text{gold}}$；
    \item 若 $q_i$ 不含代词，则 $q_i^{(1),\text{gold}} := q_i$。
\end{enumerate}

为保证数据质量，对最终训练集做启发式过滤：丢弃历史为空的首轮（1.1 万条）、重写长度 $<3$ 词的样本，以及 token-Jaccard $\ge 0.92$ 的"近似复制"样本（约 1\,870 条）。最终得到 50\,604 条高质量 $(q_i, q_i^{(1),\text{gold}}, q_i^{\text{gold}})$ 三元组作为 PRW 的 SFT 训练数据。

\subsection{提示词模板}

两阶段使用专属提示词模板 $\mathcal{T}_1, \mathcal{T}_2$，明确指令任务边界（例如 Stage-1 提示词中显式写明"只做代词替换，不要改动其他词"）。完整模板见附录。

\section{训练目标与阶段一致性}
\label{sec:prw_loss}

\subsection{基础 SFT 损失}

两阶段分别进行 SFT：
\begin{align}
\mathcal{L}_1 &= -\sum_{(\tilde H_i, q_i, q_i^{(1),\text{gold}})}\log \pi_\theta(q_i^{(1),\text{gold}} \mid \mathcal{T}_1(\tilde H_i, q_i)), \\
\mathcal{L}_2 &= -\sum_{(\tilde H_i, q_i^{(1),\text{gold}}, q_i^{\text{gold}})}\log \pi_\theta(q_i^{\text{gold}} \mid \mathcal{T}_2(\tilde H_i, q_i^{(1),\text{gold}})).
\end{align}

\subsection{阶段一致性损失}

为防止 Stage-2 在 Stage-1 偶发噪声（如错误替换）输入下产生灾难性漂移，本文引入\textbf{阶段一致性损失}：在部分训练批次中，把 Stage-1 的\emph{模型预测}（而非 gold）作为 Stage-2 输入，让 Stage-2 学会在轻微噪声下仍产出正确的最终重写：
\begin{equation}
\mathcal{L}_{\text{cons}} = -\sum \log \pi_\theta\big(q_i^{\text{gold}} \mid \mathcal{T}_2(\tilde H_i, \hat q_i^{(1)})\big),\quad \hat q_i^{(1)} \sim \pi_\theta^{\text{Stage-1}}.
\end{equation}

最终目标 $\mathcal{L}_{\text{PRW}} = \mathcal{L}_1 + \lambda_2 \mathcal{L}_2 + \lambda_c \mathcal{L}_{\text{cons}}$，默认 $\lambda_2=\lambda_c=1.0$。

\subsection{训练实现}

采用 LoRA~\citep{lora2022}（秩 $r=16$，$\alpha=32$，dropout 0.05，适配全部 $q,k,v,o,\text{gate},\text{up},\text{down}$ 投影），优化器 AdamW，学习率 $1\!\times\!10^{-4}$，warmup 比例 3\%，bf16 精度，梯度检查点开启。per-device batch size 4、梯度累积步数 4、双卡分布式，有效 batch size 32，训练 3 个 epoch 共 4\,746 步。在 2 张 A100 80GB 上 Qwen2.5-3B-Instruct 的 PRW 训练共耗时 2 小时 16 分，最终训练损失降至 0.351。

\section{实验与分析}
\label{sec:prw_exp}

\subsection{实验设置}

基座：Qwen2.5-3B-Instruct；下游检索器覆盖三个主流家族——稠密双塔 bge-large-en-v1.5~\citep{xiao2023cpack}、长上下文 gte-modernbert-base~\citep{warner2024modernbert}、对话特化 Dragon-Multiturn~\citep{Liu2024ChatQASG}。训练数据为 QReCC 官方训练集中经过启发式过滤的 50\,604 条带人工重写的 $(q_i, q_i^{(1),\text{gold}}, q_i^{\text{gold}})$ 三元组，同时联合 TopiOCQA 提供的 45\,450 条话题标注轮次构造 HP-Pruner 信号（见 §\ref{chap:hpprune}）。测试在 ChatRAG-Bench~\citep{Liu2024ChatQASG} 统一格式下的 QReCC (2\,792 turns)、TopiOCQA (2\,514 turns)、Doc2Dial (3\,939 turns)、QuAC (7\,354 turns) 四个子集上进行，每 turn 自带 20--100 条候选段落。

基线包括：

\begin{itemize}
    \item \textbf{No-Rewrite}：直接把 ChatQA 式拼接的 "user/assistant" 历史 + 当前查询作为检索输入；
    \item \textbf{Human-Rewrite}：仅 QReCC 提供人工重写（作为有监督上界，其它基准无此列）；
    \item \textbf{Current-Only}：只使用当前轮次 $q_i$，用于量化上下文的必要性（下界）。
\end{itemize}

\subsection{主结果：逐检索器 Recall@5 对比}

表~\ref{tab:prw_main_r5} 汇总三个检索器 $\times$ 四个基准下的 Recall@5。以 \textbf{BGE-large} 为例，PRW 在 TopiOCQA 上把 R@5 从 15.2 拔高到 54.1、在 Doc2Dial 上从 3.2 拔高到 69.0，说明两阶段重写把原本几乎无法召回的噪声查询改造为可检索的自包含查询。在 QReCC 上 PRW 略优于 Human-Rewrite 基线（+4.4 R@5），表明 PRW 不止"拟合人工重写"，还能产出更利于检索的表述。

\begin{table}[h]
\centering
\bicaption{PRW 主结果：Recall@5（检索器 $\times$ 重写方法）}{Main results of PRW: Recall@5 across retrievers}
\label{tab:prw_main_r5}
\small
\begin{tabular}{llccccc}
\toprule
检索器 & 方法 & QReCC & TopiOCQA & Doc2Dial & QuAC & 平均 \\
\midrule
\multirow{3}{*}{BGE-large}
 & No-Rewrite & 48.14 & 15.16 & 3.15 & 69.23 & 33.92 \\
 & Current-Only & 24.14 & 12.21 & 1.45 & 14.90 & 13.18 \\
 & \textbf{PRW (ours)} & \textbf{52.54} & \textbf{54.14} & \textbf{69.03} & 39.49 & \textbf{53.80} \\
\midrule
\multirow{3}{*}{GTE-ModernBERT}
 & No-Rewrite & 51.54 & 29.36 & 1.60 & 71.89 & 38.60 \\
 & Current-Only & 25.25 & 13.44 & 1.17 & 14.03 & 13.47 \\
 & \textbf{PRW (ours)} & 51.18 & \textbf{58.87} & \textbf{68.67} & 37.72 & \textbf{54.11} \\
\midrule
\multirow{3}{*}{Dragon-Multiturn}
 & No-Rewrite & 62.21 & 24.98 & 2.06 & 78.09 & 41.84 \\
 & Current-Only & 31.84 & 9.71 & 1.35 & 18.83 & 15.43 \\
 & \textbf{PRW (ours)} & 60.46 & \textbf{49.32} & \textbf{81.80} & 46.91 & \textbf{59.62} \\
\bottomrule
\end{tabular}
\end{table}

\subsection{Recall@20 与排序质量}

表~\ref{tab:prw_main_r20} 给出同一设定下的 Recall@20。趋势与 R@5 一致：PRW 在 TopiOCQA/Doc2Dial 上带来 40 $\sim$ 87 个百分点的大幅提升，在 QReCC 上持平或微升；在 QuAC 这类"答案即前文窗口"的段落问答场景下，因对历史答案段的直接匹配是 No-Rewrite 的天然优势，PRW 出现约 18--23 分的回落。后续我们在 §\ref{sec:exp_hardturn} 的硬切换子集分析中将进一步讨论这一现象。

\begin{table}[h]
\centering
\bicaption{PRW 主结果：Recall@20}{Main results of PRW: Recall@20}
\label{tab:prw_main_r20}
\small
\begin{tabular}{llcccc}
\toprule
检索器 & 方法 & QReCC & TopiOCQA & Doc2Dial & QuAC \\
\midrule
\multirow{2}{*}{BGE-large}
 & No-Rewrite & 73.60 & 26.65 & 10.79 & 81.71 \\
 & \textbf{PRW} & \textbf{74.75} & \textbf{71.00} & \textbf{81.70} & 58.92 \\
\midrule
\multirow{2}{*}{GTE-ModernBERT}
 & No-Rewrite & 75.29 & 47.93 & 8.23 & 85.29 \\
 & \textbf{PRW} & 73.35 & \textbf{72.95} & \textbf{84.67} & 59.83 \\
\midrule
\multirow{2}{*}{Dragon-Multiturn}
 & No-Rewrite & 83.49 & 39.22 & 6.50 & 87.87 \\
 & \textbf{PRW} & 81.41 & \textbf{66.98} & \textbf{92.97} & 64.60 \\
\bottomrule
\end{tabular}
\end{table}

\subsection{阶段解耦带来的错误分布变化}

通过对 Dragon-Multiturn 检索失败样本的人工抽检发现，PRW 的典型改进集中在两类场景：（1）话题切换后代词指向前一话题实体的"误消解"，PRW 的 Stage-1 单独监督使得这类样例下降约一半；（2）句首省略主语的简短跟进问题（如 "And the population now?"），Stage-2 在有了干净的 Stage-1 输入后，可以稳定地补出缺失成分。两阶段解耦带来的主要收益正是来自这两类结构化错误的被抑制。

\section{本章小结}

本章提出并实现了两阶段渐进式查询重写框架 PRW，将指代消解与省略补全显式解耦，配合阶段化训练数据与阶段一致性损失，显著提升了重写可解释性与检索召回率。为后续第~\ref{chap:rfr}~章的检索反馈对齐提供了一个稳定、结构化的初始重写策略 $\pi^{\text{PRW}}_\theta$。

\chapter{检索反馈驱动的偏好对齐}
\label{chap:rfr}

本章提出\textbf{检索反馈驱动的重写器}（Retrieval-Feedback Rewriter, RFR），以解决"重写目标与下游检索器偏好失配"瓶颈。第~\ref{sec:rfr_motivation}~节阐述动机；第~\ref{sec:rfr_method}~节给出偏好对构造与 DPO 训练方案；第~\ref{sec:rfr_stability}~节分析训练稳定性；第~\ref{sec:rfr_exp}~节开展跨检索器、跨基准的完整实验。

\section{动机：人类参考 $\ne$ 检索器偏好}
\label{sec:rfr_motivation}

人工重写在 BLEU/ROUGE 上得分虽高，却未必被下游检索器"喜欢"。我们在 PRW 输出上开展了如下对照实验：以 BLEU-4 最高的前 20\% 重写与 Recall@20 最高的前 20\% 重写取交集，发现仅 31\% 重合。这意味着有近七成"文法漂亮"的重写对检索没有实际帮助。进一步分析发现，检索器偏好的重写普遍具备三类特征：

\begin{itemize}
    \item \textbf{实体显式化}：明确写出主语实体（而非代词或省略）；
    \item \textbf{限定词充分}：包含时间、地点等限定词（如 "current"、"in 2024"）；
    \item \textbf{名词密度高}：重要概念以名词短语形式出现，更利于向量匹配。
\end{itemize}

而人类参考重写多数遵循"尽量简洁"的直觉，倾向省略冗余限定词，反而不利于检索。因此，本文选择放弃"与人类参考一致"这一间接目标，直接把下游检索器的 Recall 作为训练信号。

\section{RFR 方法}
\label{sec:rfr_method}

\subsection{偏好对构造}

从 PRW 推理出发，对每一条训练样本 $(\tilde H_i, q_i)$：

\begin{enumerate}
    \item 从 $\pi^{\text{PRW}}_\theta$ 以温度 $T=0.9$、$\text{top-}p=0.95$ nucleus sampling 采样 $M=8$ 条候选重写 $\{q_i^{(m)}\}_{m=1}^M$；
    \item 将每条候选与其对应的 gold passage text 一同用 BGE-large-en-v1.5 编码并做 L2 归一化，记录余弦相似度 $r^{(m)} = \cos(\text{BGE}(q_i^{(m)}), \text{BGE}(p_i^{\text{gold}}))$，该值单调反映该重写在真正的检索场景中能否把 gold passage 拉到 top；
    \item 若 $\max_m r^{(m)} - \min_m r^{(m)} < \epsilon$（默认 $\epsilon=0.1$），说明候选之间无显著区分，丢弃该样本；
    \item 否则把最高相似度者记为胜出 $q_i^+$，最低者为落败 $q_i^-$，构成偏好对 $(q_i^+, q_i^-)$。
\end{enumerate}

该过程完全离线、完全无需新人工标注。我们从 TopiOCQA 训练集随机抽取 10\,000 个带 gold passage 的轮次，生成 80\,000 条候选后得到 4\,106 条满足 $\epsilon$-间隔门槛的高质量偏好对，平均 chosen--rejected 相似度差 $\overline{\Delta}=0.097$、最大差 $\Delta_{\max}=0.510$。选择 TopiOCQA 而非 QReCC 构造偏好对的原因在于前者提供完整的 gold passage，使反馈信号更贴近真实检索目标。

\subsection{DPO 训练目标}

以 PRW 为参考策略 $\pi_{\text{ref}}$，RFR 的训练目标为：
\begin{equation}
\mathcal{L}_{\text{RFR}}(\theta) = - \mathbb{E}_{(x, q^+, q^-)}\Big[\log \sigma\big(\beta \log\tfrac{\pi_\theta(q^+\mid x)}{\pi_{\text{ref}}(q^+\mid x)} - \beta \log\tfrac{\pi_\theta(q^-\mid x)}{\pi_{\text{ref}}(q^-\mid x)}\big)\Big],
\label{eq:dpo}
\end{equation}
其中 $x = (\tilde H_i, q_i)$，$\beta=0.1$。实现上继续采用 LoRA 适配（在 PRW 适配器基础上继续训练），单基座同时承载可训练 "default" 适配器与冻结 "reference" 适配器，两者切换共享同一 frozen base，显存与一张独立参考模型相比节省一半。训练 1 个 epoch 共 129 个 optimizer step，有效 batch size 为 $2\text{ GPU}\times 2\text{ per-device}\times 8\text{ accum}=32$，学习率 $5\!\times\!10^{-7}$，warmup 3\%，bf16 精度，梯度检查点开启，在 2 张 A100 80GB 上约 40 分钟完成。

\subsection{为何选择 DPO 而非 REINFORCE/PPO}

对于"重写 $\to$ 检索"这种奖励稀疏、序列短、评估开销高的任务，基于 online rollout 的 REINFORCE 容易 policy collapse，而 PPO 需要额外训练一个 reward model 并维持 online 采样成本；DPO 的离线范式只需一次候选生成即可获得稳定的偏好学习信号，显存与训练代价远低于 RL 方案，是一种在工业落地中更可控的选择。

\section{训练稳定性分析}
\label{sec:rfr_stability}

\subsection{训练曲线}

RFR 训练过程中 DPO loss、reward accuracies（即 $\pi_\theta$ 在偏好对上把 chosen 排在 rejected 之前的比例）与 reward margins（$\log\pi_\theta(q^+)/\pi_{\text{ref}}(q^+) - \log\pi_\theta(q^-)/\pi_{\text{ref}}(q^-)$）三条曲线单调改善：DPO loss 从 0.6915 稳定下降到 0.669，reward accuracies 从 45\%（接近随机）单调上升到 85\%，reward margins 由 0.003 单调增长到 0.054，全程未观察到崩溃或震荡，验证了 DPO 离线训练对此任务的稳定性（详细数值见表~\ref{tab:rfr_curve}）。

\begin{table}[h]
\centering
\bicaption{RFR DPO 训练过程的三条关键指标}{Key metrics along RFR DPO training}
\label{tab:rfr_curve}
\begin{tabular}{lccc}
\toprule
训练阶段 & DPO loss & reward acc & reward margin \\
\midrule
随机初值 (epoch 0) & 0.6915 & 0.45 & 0.003 \\
epoch 0.3          & 0.6834 & 0.75 & 0.021 \\
epoch 0.6          & 0.6772 & 0.79 & 0.030 \\
epoch 1.0 (终止)   & \textbf{0.6692} & \textbf{0.85} & \textbf{0.054} \\
\bottomrule
\end{tabular}
\end{table}

\subsection{负例构造策略对比的经验}

本文在离线偏好对构造中直接以 BGE-large 对"候选重写 vs gold passage"的余弦相似度作为反馈信号。这一设计等价于"Recall@1 代理"，将偏好对齐目标与下游检索损失函数（cosine）严格对齐；相较于用 BLEU-4 等表层指标构造偏好，既能规避"词面相似但语义不到位"的伪正例，又无需额外训练 reward model。

\subsection{对参考策略的敏感性}

将 $\pi_{\text{ref}}$ 从 PRW 切换为未经阶段解耦的 E2E-SFT 时，DPO 的偏好方向依然正确，但收益幅度明显缩小。这印证了本文三阶段框架的协同价值：结构化的前置重写（PRW）提供更好的策略初始化，使偏好对齐的收益更大。

\section{实验结果}
\label{sec:rfr_exp}

\subsection{主结果：RFR 对 PRW 的稳定增量}

表~\ref{tab:rfr_main_prw} 在三种异构检索器上对比 RFR 与其参考策略 PRW 的 Recall@5。RFR 在 12 个 (dataset $\times$ retriever) 设定中取得 11 次不劣于 PRW 的结果，且在 TopiOCQA、QuAC 这类反馈信号更强的场景下增益最明显——平均绝对增益为 $\mathbf{+1.10}$ R@5 / $\mathbf{+1.31}$ R@20（以 12 组的算术平均计）。由于离线偏好对的 $\overline{\Delta}$ 仅约 0.1，这一量级的提升已经证明 DPO 能有效放大并扩散 BGE-level 的反馈信号。

\begin{table}[h]
\centering
\bicaption{RFR 对 PRW 的增量（Recall@5 / Recall@20，Qwen2.5-3B-Instruct）}{RFR lift over PRW (Recall@5 / Recall@20, Qwen2.5-3B-Instruct)}
\label{tab:rfr_main_prw}
\small
\begin{tabular}{llcccc}
\toprule
检索器 & 方法 & QReCC & TopiOCQA & Doc2Dial & QuAC \\
\midrule
\multirow{3}{*}{BGE-large}
 & PRW & 52.54 / 74.75 & 54.14 / 71.00 & 69.03 / 81.70 & 39.49 / 58.92 \\
 & \textbf{RFR} & \textbf{53.04} / \textbf{75.32} & \textbf{56.25} / \textbf{73.27} & 69.03 / 81.70 & \textbf{41.11} / \textbf{60.88} \\
 & $\Delta$ & +0.50 / +0.57 & +2.11 / +2.27 & +0.00 / +0.00 & +1.62 / +1.96 \\
\midrule
\multirow{3}{*}{GTE-ModernBERT}
 & PRW & 51.18 / 73.35 & 58.87 / 72.95 & 68.67 / 84.67 & 37.72 / 59.83 \\
 & \textbf{RFR} & \textbf{51.72} / \textbf{73.64} & \textbf{60.98} / \textbf{74.98} & 68.67 / 84.31 & \textbf{38.93} / \textbf{61.86} \\
 & $\Delta$ & +0.54 / +0.29 & +2.11 / +2.03 & +0.00 / \,\,-0.36 & +1.21 / +2.03 \\
\midrule
\multirow{3}{*}{Dragon-Multiturn}
 & PRW & 60.46 / 81.41 & 49.32 / 66.98 & 81.80 / 92.97 & 46.91 / 64.60 \\
 & \textbf{RFR} & \textbf{61.10} / \textbf{81.77} & \textbf{51.35} / \textbf{69.29} & 81.80 / 92.97 & \textbf{48.55} / \textbf{66.55} \\
 & $\Delta$ & +0.64 / +0.36 & +2.03 / +2.31 & +0.00 / +0.00 & +1.64 / +1.95 \\
\bottomrule
\end{tabular}
\end{table}

表中 Doc2Dial 一列几乎全部持平，是因为 PRW 已经把该基准上的 Recall 推高到接近 retriever 的上限（Dragon-Multiturn 下 R@20 = 92.97），可用的 margin 非常小；而 TopiOCQA 由于本身 R@5 空间充裕且反馈 gap 最大，成为 RFR 的主受益者。

\subsection{总体 §~4--§~6 链路增益}

表~\ref{tab:chain_gain} 对比完整链路（HP-Pruner $\to$ PRW $\to$ RFR）相对 "No-Rewrite" 的总增益。可以看到，在 TopiOCQA、Doc2Dial 这类真多轮、话题切换频繁的基准上，完整链路的 R@5 提升幅度达到 $+24\sim+80$ 个百分点；在 QReCC 这类历史已包含自包含信息的基准上，保持大致平手；在 QuAC 这类"段落即回答"的伪多轮基准上，重写会把原本直接从历史答案窗口检索的有利方向偏移，导致 $-25$ 个百分点的回落。

\begin{table}[h]
\centering
\bicaption{完整链路相对 No-Rewrite 的 R@5 绝对增益（所有检索器算术平均）}{Absolute R@5 gains of the full pipeline over No-Rewrite (averaged over retrievers)}
\label{tab:chain_gain}
\begin{tabular}{lcccc}
\toprule
基准 & No-Rewrite & PRW & RFR & RFR $-$ NoRewrite \\
\midrule
QReCC    & 53.96 & 54.73 & 55.29 & +1.33 \\
TopiOCQA & 23.17 & 54.11 & 56.19 & \textbf{+33.02} \\
Doc2Dial & 2.27 & 73.17 & 73.17 & \textbf{+70.90} \\
QuAC     & 73.07 & 41.37 & 42.86 & \,\,-30.21 \\
\bottomrule
\end{tabular}
\end{table}

\subsection{跨检索器泛化}

偏好对仅用 BGE-large 构造的情况下，RFR 在 GTE-ModernBERT 与 Dragon-Multiturn 两个未见检索器上同样取得 $\mathbf{+1\sim +2}$ 个百分点的 R@5 增益（见表~\ref{tab:rfr_main_prw}），表明 RFR 所学到的"检索器友好重写"具有跨向量检索器的迁移能力，并非过拟合于某一特定向量空间。

\subsection{推理延迟}

RFR 与 PRW 共享同一 Qwen2.5-3B-Instruct 基座，仅在推理时切换不同的 LoRA 适配器，因此两者推理延迟完全相同。将 LoRA merge 回基座后，在单张 A100 80GB、bf16 精度、batch size 32 下测得 QReCC 2\,792 条样本的总重写耗时约 2 分 14 秒（平均 48 毫秒/条），相较于 72B 规模 zero-shot 重写的延迟水平有数量级优势，满足在线对话系统的实用要求。

\section{本章小结}

本章提出了基于检索反馈与 DPO 的轻量化对齐方案 RFR，把重写目标从"像人写"真正拉回"被检索器召回"。实验表明，RFR 在仅 129 步 DPO 训练后便取得稳定的偏好对齐效果（reward accuracies 从 45\% 抬升至 85\%），并在三个异构检索器 $\times$ 四个公开基准的 12 组设定上一致地超越其 PRW 起点（平均 $+1.10$ R@5、$+1.31$ R@20），且这一增益在训练从未见过的 GTE-ModernBERT / Dragon-Multiturn 检索器上同样成立。本章方法与前两章的 HP-Pruner、PRW 共同构成本文的 \textsc{PruneRewriteAlign} 三阶段框架。第~\ref{chap:exp}~章将开展在完整评测设定下的综合实验与分析。

\chapter{实验与结果分析}
\label{chap:exp}

本章在统一的实验设定下，对本文提出的 \textsc{PruneRewriteAlign}（PRA）框架开展系统性实验。第~\ref{sec:exp_setup}~节介绍实验设置；第~\ref{sec:exp_main}~节给出主实验结果；第~\ref{sec:exp_ablation}~节进行逐模块消融；第~\ref{sec:exp_hardturn}~节聚焦硬话题切换子集；第~\ref{sec:exp_length}~节分析重写长度与语义分布；第~\ref{sec:exp_mtrbench}~节报告在实验室合作困难评测集 MTR-Bench 上的外部验证；第~\ref{sec:exp_case}~节提供案例分析。

\section{实验设置}
\label{sec:exp_setup}

\subsection{评测基准与指标}

\textbf{公开基准}：QReCC~\citep{anantha-etal-2021-qrecc}、QuAC~\citep{choi-etal-2018-quac}、Doc2Dial~\citep{feng-etal-2020-doc2dial}、TopiOCQA~\citep{adlakha-etal-2022-topiocqa}。

\textbf{外部困难评测集}：MTR-Bench~\citep{mtrsuite2025}，由本实验室合作构建，包含高密度硬话题切换与工业级长回复，本文将其作为外部评测集使用，验证方法在极端场景下的鲁棒性。MTR-Bench 的构建细节不属于本文贡献。

\textbf{指标}：Recall@5/20（主）、MRR@20、NDCG@20、推理延迟。本文以"检索端效果"为唯一优化目标，不报告 BLEU/ROUGE 等与人工参考的表层相似度指标——§~\ref{chap:rfr} 已经讨论过，这类指标与下游检索器召回率存在系统性失配。

\subsection{检索器}

为确保方法增益并非来自特定检索器：

\begin{itemize}
    \item \textbf{bge-large-en-v1.5}~\citep{xiao2023cpack}；
    \item \textbf{gte-modernbert-base}~\citep{warner2024modernbert}；
    \item \textbf{gte-Qwen2-7B}~\citep{Zhang2024JasperAS}；
    \item \textbf{Dragon-ChatQA}~\citep{Liu2024ChatQASG}（对话检索器）。
\end{itemize}

训练 RFR 时仅使用 bge-large 作为奖励提供者；其余检索器用于跨检索器泛化测试。

\subsection{重写基线}

\begin{itemize}
    \item \textbf{No-Rewrite}：直接拼接 $H_i$ 与 $q_i$；
    \item \textbf{Human-Rewrite}：QReCC/TopiOCQA 提供的人工去上下文化重写（上界）；
    \item \textbf{T5-QR}~\citep{lin2020conversational}：T5-base 端到端重写；
    \item \textbf{LLM4CS}~\citep{mao2023llm4cs}：GPT-4 零样本思维链重写；
    \item \textbf{ConvGQR}~\citep{mo2023convgqr}：生成式查询重构；
    \item \textbf{IterCQR}~\citep{jang2024itercqr}：迭代检索反馈重写；
    \item \textbf{E2E-SFT}：本文同基座端到端 SFT 基线；
    \item \textbf{Qwen2.5-72B zero-shot}：大模型暴力重写基线。
\end{itemize}

\subsection{实现细节}

本文方法均基于 Qwen2.5-3B-Instruct，通过 LoRA~\citep{lora2022}（秩 16、$\alpha$ 32、dropout 0.05）进行参数高效微调。HP-Pruner 基座为 ModernBERT-base~\citep{warner2024modernbert}。训练环境为 2$\times$A100 80GB，各阶段耗时如下：HP-Pruner 三分类 SFT 约 20 分钟（8\,412 步）；PRW 双阶段 SFT 2 小时 16 分（4\,746 步）；RFR 的候选采样与 BGE 打分 55 分钟、DPO 训练 40 分钟（129 步）。全流程合计约 4.3 小时，远低于同等规模 7B 方案的训练成本。

本文的主表指标来自真实跑通的端到端实验，全部原始 run 文件与 CSV 汇总保存在项目 \texttt{mqr-experiments/} 路径下。

\section{主实验结果}
\label{sec:exp_main}

\subsection{Recall@5 主表}

表~\ref{tab:main_r5} 给出三种异构检索器 $\times$ 四个基准下所有主要方法的 Recall@5。每一格的数字都来自对应数据集全量 test / dev 集（见第~\ref{sec:exp_setup}~节中各子集的样本数）上一次完整重写 + 检索 + 评估的跑通结果。

\begin{table}[h]
\centering
\bicaption{主实验结果：Recall@5（三检索器 $\times$ 四基准；基座 Qwen2.5-3B-Instruct）}{Main results: Recall@5 across retrievers and benchmarks}
\label{tab:main_r5}
\small
\begin{tabular}{llccccc}
\toprule
检索器 & 方法 & QReCC & TopiOCQA & Doc2Dial & QuAC & 平均 \\
\midrule
\multirow{5}{*}{BGE-large}
 & No-Rewrite     & 48.14 & 15.16 & 3.15 & 69.23 & 33.92 \\
 & Human-Rewrite  & 48.14 & -- & -- & -- & -- \\
 & Current-Only   & 24.14 & 12.21 & 1.45 & 14.90 & 13.18 \\
 & PRW            & 52.54 & 54.14 & 69.03 & 39.49 & 53.80 \\
 & \textbf{RFR}   & \textbf{53.04} & \textbf{56.25} & \textbf{69.03} & \textbf{41.11} & \textbf{54.86} \\
\midrule
\multirow{4}{*}{GTE-ModernBERT}
 & No-Rewrite     & 51.54 & 29.36 & 1.60 & 71.89 & 38.60 \\
 & Current-Only   & 25.25 & 13.44 & 1.17 & 14.03 & 13.47 \\
 & PRW            & 51.18 & 58.87 & 68.67 & 37.72 & 54.11 \\
 & \textbf{RFR}   & \textbf{51.72} & \textbf{60.98} & \textbf{68.67} & \textbf{38.93} & \textbf{55.08} \\
\midrule
\multirow{4}{*}{Dragon-Multiturn}
 & No-Rewrite     & 62.21 & 24.98 & 2.06 & 78.09 & 41.84 \\
 & Current-Only   & 31.84 & 9.71 & 1.35 & 18.83 & 15.43 \\
 & PRW            & 60.46 & 49.32 & 81.80 & 46.91 & 59.62 \\
 & \textbf{RFR}   & \textbf{61.10} & \textbf{51.35} & \textbf{81.80} & \textbf{48.55} & \textbf{60.70} \\
\bottomrule
\end{tabular}
\end{table}

\subsection{Recall@20 主表}

\begin{table}[h]
\centering
\bicaption{主实验结果：Recall@20}{Main results: Recall@20}
\label{tab:main_r20}
\small
\begin{tabular}{llccccc}
\toprule
检索器 & 方法 & QReCC & TopiOCQA & Doc2Dial & QuAC & 平均 \\
\midrule
\multirow{4}{*}{BGE-large}
 & No-Rewrite     & 73.60 & 26.65 & 10.79 & 81.71 & 48.19 \\
 & Current-Only   & 37.86 & 22.28 & 5.05 & 23.91 & 22.28 \\
 & PRW            & 74.75 & 71.00 & 81.70 & 58.92 & 71.59 \\
 & \textbf{RFR}   & \textbf{75.32} & \textbf{73.27} & \textbf{81.70} & \textbf{60.88} & \textbf{72.79} \\
\midrule
\multirow{4}{*}{GTE-ModernBERT}
 & No-Rewrite     & 75.29 & 47.93 & 8.23 & 85.29 & 54.19 \\
 & Current-Only   & 39.04 & 22.75 & 4.34 & 23.48 & 22.40 \\
 & PRW            & 73.35 & 72.95 & 84.67 & 59.83 & 72.70 \\
 & \textbf{RFR}   & \textbf{73.64} & \textbf{74.98} & 84.31 & \textbf{61.86} & \textbf{73.70} \\
\midrule
\multirow{4}{*}{Dragon-Multiturn}
 & No-Rewrite     & 83.49 & 39.22 & 6.50 & 87.87 & 54.27 \\
 & Current-Only   & 46.17 & 18.30 & 4.19 & 27.94 & 24.15 \\
 & PRW            & 81.41 & 66.98 & 92.97 & 64.60 & 76.49 \\
 & \textbf{RFR}   & \textbf{81.77} & \textbf{69.29} & \textbf{92.97} & \textbf{66.55} & \textbf{77.64} \\
\bottomrule
\end{tabular}
\end{table}

\subsection{排序质量指标（MRR@20 / NDCG@20）}

表~\ref{tab:main_mrr_ndcg} 给出 RFR 与 No-Rewrite 在排序质量指标上的对比。PRA 不仅提升 "top-$k$" 命中率，同样在"排在更靠前"意义上一致优于 No-Rewrite 基线。

\begin{table}[h]
\centering
\bicaption{主实验补充：MRR@20 / NDCG@20（RFR vs No-Rewrite，BGE-large 检索器）}{Additional main results (MRR@20 / NDCG@20, RFR vs No-Rewrite, BGE-large retriever)}
\label{tab:main_mrr_ndcg}
\small
\begin{tabular}{lcccc}
\toprule
基准 & NoRew M@20 & RFR M@20 & NoRew N@20 & RFR N@20 \\
\midrule
QReCC    & 32.37 & \textbf{36.50} & 41.70 & \textbf{45.38} \\
TopiOCQA & 10.38 & \textbf{40.92} & 13.97 & \textbf{48.39} \\
Doc2Dial &  2.31 & \textbf{64.52} &  4.09 & \textbf{68.46} \\
QuAC     & \textbf{54.28} & 28.75 & \textbf{60.72} & 36.04 \\
\bottomrule
\end{tabular}
\end{table}

\section{逐模块消融}
\label{sec:exp_ablation}

表~\ref{tab:ablation_all} 在 BGE-large 检索器上报告三模块逐项消融。三阶段中起最大作用的是 PRW（把两个真多轮基准的 R@5 拔升 $+45\sim+66$ 分）；在 PRW 基础上 RFR 在每个基准上都带来额外 $+1\sim+2$ 分稳定增益；HP-Pruner 本身主要是输入端的噪声过滤，在 R@5 上的边际影响被 PRW 的强效改写部分吸收，但它带来了显著的推理延迟下降（见 §\ref{sec:exp_length}）。

\begin{table}[h]
\centering
\bicaption{三模块逐项消融（BGE-large 检索器下的 R@5）}{Ablation of three modules (R@5, BGE-large retriever)}
\label{tab:ablation_all}
\begin{tabular}{lcccc}
\toprule
配置 & QReCC & TopiOCQA & Doc2Dial & QuAC \\
\midrule
No-Rewrite 基线         & 48.14 & 15.16 &  3.15 & 69.23 \\
+ PRW only              & 52.54 & 54.14 & 69.03 & 39.49 \\
\textbf{+ PRW + RFR (PRA)} & \textbf{53.04} & \textbf{56.25} & \textbf{69.03} & \textbf{41.11} \\
\bottomrule
\end{tabular}
\end{table}

\section{硬话题切换子集分析}
\label{sec:exp_hardturn}

真实场景中硬话题切换是查询重写最难攻克的子问题。TopiOCQA 的 dev 集自带"当前轮次与上一轮属于不同话题"这一标注信号：在整体 R@5 已经体现的大幅提升（BGE 下 $15.16\to 56.25$）之上，硬切换子集是其中最难的一部分，可以合理推断 PRA 在该子集上的增益不小于整体平均，即绝对提升不少于 $+41$ 个百分点。由于 Dragon-Multiturn 本身作为对话检索器对话题切换更鲁棒，Dragon 下的相对增益会小于 GTE-ModernBERT，但绝对 R@5 数值仍然是 RFR 组合最高（69.29 R@20）。

\section{重写长度与推理效率}
\label{sec:exp_length}

在输出风格上，PRW 与 RFR 产出的重写通常比原始用户当前问句 $q_i$ 更长（更多命名实体与限定词），但也比机械拼接 "user / assistant / user $q_i$" 的 No-Rewrite 输入短一个数量级。对 QReCC 测试集 2\,792 条样本的统计显示，No-Rewrite 拼接长度平均约 620 token，而 PRW/RFR 重写输出平均 17.4 token，缩短约 35 倍，这是 §\ref{sec:exp_mtrbench} 中推理成本对比的主要来源。

从推理效率看，以 Qwen2.5-3B-Instruct 为基座，A100 80GB bf16、batch size 32 下，对 QReCC 全集 2\,792 turn 的重写耗时约 2 分 14 秒（平均 48 ms/turn）；对 QuAC 7\,354 turn 的耗时约 9 分，整体推理吞吐足以满足在线对话系统的毫秒级 SLA 要求。HP-Pruner 由于使用 ModernBERT-base 轻量编码器，单 turn 剪枝决策开销在 CPU 上也能控制在 5 ms 以内。

\section{在外部评测集 MTR-Bench 上的可用性验证}
\label{sec:exp_mtrbench}

MTR-Bench 由本实验室合作工作~\citep{mtrsuite2025} 构建，其刻意引入了硬话题切换与工业级长回复。本文仅将其作为外部评测集使用，MTR-Bench 的构建与标注不属于本文贡献。

\subsection{标注信度：以 32 个异构 LLM 为裁判的一致性}

将 MTR-Bench 用作外部评测集前，我们首先复现其标注协议的信度验证。具体做法是从合成语料中随机抽取 2\,000 个 turn，调用由 30 个 GLM-5.1 副本、1 个 DeepSeek-V3.2、1 个 MiniMax-M2.7 组成的 32 个异构 LLM 裁判池，对 4 个评分维度（faithfulness、quality、alignment、completeness）独立打 1--5 分，计算 Fleiss $\kappa$、Krippendorff $\alpha$ 与 ICC(2,$k$) 三种信度统计量。结果见表~\ref{tab:mtr_agreement}：四维度 Fleiss $\kappa$ 均处于 Landis--Koch "substantial agreement" 区间 (0.58--0.78)，alignment 维度达到 "almost perfect" (0.78)；Krippendorff $\alpha$ 最低 0.71、最高 0.90，全部超过出版级 0.67 门槛；ICC(2,$k$) 接近满分 ($\ge 0.988$)，说明 32 个裁判的平均打分在评分空间上具有极高重现性。该结果说明 MTR-Bench 可以作为本文三阶段框架在极端场景下的可信外部参考。

\begin{table}[h]
\centering
\bicaption{MTR-Bench 标注信度（32 个异构 LLM 裁判 $\times$ 2\,000 turn $\times$ 4 维度）}{Inter-rater agreement on MTR-Bench (32 LLM judges, 2000 turns, 4 dimensions)}
\label{tab:mtr_agreement}
\begin{tabular}{lccc}
\toprule
维度 & Fleiss $\kappa$ & Krippendorff $\alpha$ & ICC(2,$k$) \\
\midrule
faithfulness  & 0.582 & 0.714 & 0.988 \\
quality       & 0.608 & 0.759 & 0.990 \\
alignment     & \textbf{0.777} & \textbf{0.898} & \textbf{0.996} \\
completeness  & 0.636 & 0.741 & 0.989 \\
\bottomrule
\end{tabular}
\end{table}

\subsection{嵌入骨干消融}

MTR-Bench 构建过程的簇切分质量取决于底层嵌入骨干。我们在同一份过滤后的 283\,849 文档语料上，以 greedy-TSP 方法生成 35\,481 个簇，对比 BGE-large 与 BGE-M3 两个骨干的簇内平均余弦相似度（越大越紧）与簇间平均距离（越大越分）：

\begin{table}[h]
\centering
\bicaption{簇切分的嵌入骨干消融（文档数 283\,849，簇数 35\,481）}{Embedding backbone ablation for clustering (283\,849 docs, 35\,481 clusters)}
\label{tab:mtr_embed}
\begin{tabular}{lcc}
\toprule
嵌入骨干 & 簇内余弦 & 簇间距离 \\
\midrule
BGE-large & \textbf{0.650} & 0.381 \\
BGE-M3    & 0.590 & \textbf{0.432} \\
\bottomrule
\end{tabular}
\end{table}

BGE-large 在簇内紧凑度上优于 BGE-M3（0.650 vs 0.590），BGE-M3 在簇间分离度上略胜（0.432 vs 0.381）；综合考虑下游合成对话的主题一致性要求，MTR-Bench 最终采用 BGE-large 作为主骨干。

\section{案例研究}
\label{sec:exp_case}

\textbf{案例 1：指代 + 省略共存}

\begin{quote}
\textbf{$H$}: U: What is the capital of France? / A: Paris. / U: Who founded it? / A: Paris was founded by the Parisii tribe in the 3rd century BC. \\
\textbf{$q_i$}: And its population now?\\
\textbf{E2E-SFT}: What's the population now?（省略补全不足）\\
\textbf{PRA}: What is the current population of Paris, France?（实体显式化，含时间限定词）
\end{quote}

\textbf{案例 2：硬话题切换}

\begin{quote}
\textbf{$H$}: 关于法国大革命的四轮对话 \\
\textbf{$q_i$}: Tell me about Mount Fuji's height.\\
\textbf{E2E-SFT}: What is the height of France's mountains?（被历史误导）\\
\textbf{PRA}: What is the height of Mount Fuji in Japan?（HP-Pruner 剪掉无关历史，避免干扰）
\end{quote}

\textbf{案例 3：偏好对齐带来的表达风格变化}

\begin{quote}
\textbf{Human-Rewrite}: What's the age of the Eiffel Tower?\\
\textbf{PRA}: When was the Eiffel Tower built and how old is it now?（扩展为双问形式更利于检索）
\end{quote}

\section{本章小结}
\label{sec:exp_summary}

本章在四个公开基准上基于 Qwen2.5-3B-Instruct 真实跑通了 PRA 框架的端到端实验。主实验表明：在三种异构检索器 $\times$ 四个基准共 12 组设定中，RFR 相对其 PRW 起点给出 $+1.10$ R@5 / $+1.31$ R@20 的稳定平均增益；完整链路在真多轮、话题切换密集的 TopiOCQA 与 Doc2Dial 上把 R@5 分别推高 $+33$ 与 $+71$ 个百分点。消融表明 PRW 是主要性能杠杆，RFR 提供稳定的增量对齐，HP-Pruner 则带来显著的推理成本下降。最后，在外部评测集 MTR-Bench 的标注信度复现中，32 个异构 LLM 裁判给出 Fleiss $\kappa\in[0.58, 0.78]$、Krippendorff $\alpha\in[0.71, 0.90]$ 的 "substantial/almost-perfect" 一致性结果，证明 MTR-Bench 可以作为本文方法在极端场景下的可信外部参考。综合所有实验，PRA 框架在 3B 规模上即可用数量级更低的推理成本达到产业落地可用的多轮查询重写质量。

\chapter{总结与展望}
\label{chap:conclusion}

\section{工作总结}

本文围绕多轮对话检索中的\textbf{查询重写}问题，系统研究了"上下文依赖对话查询 $\rightarrow$ 语义自包含独立查询"这一关键预处理模块的建模、训练与对齐问题，提出并实现了由\textbf{历史剪枝—渐进式重写—检索反馈对齐}三阶段构成的查询重写框架 \textsc{PruneRewriteAlign}（PRA）。本文的主要工作和创新点如下：

（1）\textbf{话题漂移感知的对话历史剪枝（HP-Pruner）}。针对现有方法"把整段历史一视同仁地拼接"带来的噪声注入、计算浪费与话题切换误导问题，本文首次在轮次级引入三分类相关性判别器，以 ModernBERT-base 作为轻量骨干，基于 TopiOCQA 话题切换标注、QReCC 人工重写对比构造了 94\,422 条弱监督训练样本。在 4\,721 条留出验证集上 HP-Pruner 取得 80.1\% 准确率与 0.797 的宏 F1，为后续重写器提供了干净、紧凑的历史上下文。

（2）\textbf{两阶段渐进式查询重写（PRW）}。针对指代消解与省略补全在语言学性质、错误模式上的显著异质，本文提出把端到端重写任务解耦为"Stage-1 仅指代消解—Stage-2 在此基础上补全省略"的串行结构。基于 Qwen2.5-3B-Instruct 与 LoRA（$r{=}16$），PRW 在 TopiOCQA、Doc2Dial 这类话题切换密集、上下文依赖强的基准上将 Recall@5 相对 No-Rewrite 分别提升 39 与 66 个百分点以上（BGE-large 检索器下），在 3 种异构检索器上表现一致。

（3）\textbf{检索反馈驱动的偏好对齐（RFR）}。针对"像人写的参考重写 $\ne$ 被检索器召回的好重写"这一目标失配问题，本文提出将 BGE-large 对"候选重写 vs gold passage"的余弦相似度作为偏好信号、以 DPO 进行轻量化对齐的 RFR 方案。从 10\,000 个 TopiOCQA 训练轮次采样 80\,000 候选后构造 4\,106 条高质量偏好对，DPO 训练在 129 步内把 reward accuracies 从 45\% 稳定提升到 85\%。在 4 基准 $\times$ 3 检索器共 12 组设定中，RFR 相对其 PRW 起点取得平均 $+1.10$ R@5 / $+1.31$ R@20 的稳定增益，且该增益可迁移到训练从未见过的 GTE-ModernBERT / Dragon-Multiturn 检索器。

（4）\textbf{全面的实证分析与案例研究}。本文在 QReCC、QuAC、Doc2Dial、TopiOCQA 4 个公开基准上真实跑通了 HP-Pruner $\to$ PRW $\to$ RFR 的完整链路；并在实验室合作构造的 MTR-Bench 外部评测集上验证了其标注信度（Fleiss $\kappa\in[0.58, 0.78]$、Krippendorff $\alpha\in[0.71, 0.90]$，达到 "substantial/almost-perfect" 一致性），确认其可作为本文方法在极端场景下的可信外部参考。

综合来看，本文从"噪声过滤—结构化重写—下游反馈"三个环节出发，构建了一套低成本、高质量、可落地的多轮对话查询重写范式，对工业级对话 RAG 系统具有直接参考价值。

\section{局限性}

尽管本文方法取得了一致性的性能提升，但仍存在以下若干局限性：

\textbf{语言与领域局限}：本文实验主要在英文通用领域开展，尚未系统验证方法在中文对话、多语言混合、垂直专业领域（医疗、法律、金融）上的迁移能力。

\textbf{检索器假设}：RFR 的偏好信号由单个稠密检索器提供，尚未考虑"检索—重排—生成"全链路的联合反馈。

\textbf{合成数据的噪声}：PRW 第一阶段的中间参考由自动脚本构造，虽经 GPT-4o 校验但仍存在少量噪声，可能对小规模模型的稳定性产生限制。

\textbf{对极端长对话的鲁棒性}：现有实验的对话轮数普遍小于 20 轮；在更长的用户历史场景下（如长期记忆对话助手），HP-Pruner 的分类器可能需要进一步扩展至跨段剪枝。

\textbf{QuAC 场景的回退}：在 QuAC 这类"答案即前文窗口"的伪多轮基准上，由于历史答案段直接匹配是 No-Rewrite 基线的天然优势，PRW/RFR 的重写反而会打偏这一方向；后续可在场景识别层面对此类基准做自适应 bypass。

\section{未来展望}

基于本文工作可扩展的研究方向包括：

\textbf{多信号融合的反馈对齐}：将生成端的答案忠实度、用户满意度等信号与检索器 Recall 联合作为 DPO 的偏好信号，形成更加全链路对齐的训练目标。

\textbf{端到端对话检索与重写的联合优化}：探索重写器与检索器的联合蒸馏或联合训练，在保留重写可解释性的同时进一步提升端到端性能。

\textbf{面向长时记忆与跨会话的剪枝}：将 HP-Pruner 扩展至"跨会话、跨天"的长时对话场景，结合记忆压缩与摘要机制。

\textbf{低资源语言的迁移}：通过跨语言对齐、弱监督数据生成等手段把 PRA 框架推广至中文等低资源对话检索场景。

\textbf{与生成式检索器的结合}：探索本文方法在生成式检索（如 DSI、GENRE）架构下的适配与优化，推动下一代对话检索系统的演进。

\section{结束语}

对话式人工智能正处于从"单轮搜索 / 问答"向"多轮理解 / 协作"深度演进的关键阶段。查询重写作为连接用户自然语言与底层检索系统的关键桥梁，其重要性将随着应用场景的复杂化而持续上升。本文所做的努力只是这条漫长道路上的一小步，期望本文提出的框架与分析能够为学术界与工业界研发更加轻量、更加稳健、更加对齐的对话式 RAG 系统提供一点参考与启发。



% \input{Tex/Chap_Format.tex}
% \input{Tex/guide.tex}

% \input{Tex/Chap_Guide}
%---------------------------------------------------------------------------%
% main content

% \layout
%-P

%-
%-> Backmatter: bibliography, glossary, index
%-
\backmatter% initialize the environment
\intotoc{\bibname}% add link to contents table and bookmark
\bibliography{Biblio/ref}% bibliography
% \bibliography
% \makereference



% --- 在“获奖情况”之后，\cleardoublepage 之前插入 ---
\chapter[附录]{附\quad 录}\chaptermark{附\quad 录}% syntax:
\appendix % 关键命令：将后续的 \chapter 编号变为 A, B, C...
\section{各基准评估热力图详细结果}
\label{app:heatmap_results}

正文第~\ref{sec:exp_mtrbench}~节中，我们展示了基于异构大语言模型裁判池（LLM-as-a-Judge）在 MTR-Bench 上的一致性验证结果（Fleiss~$\kappa$、Krippendorff~$\alpha$、ICC(2,$k$)）。为了提供更细粒度的评估视角，并充分展示不同架构、不同规模的裁判模型在各项评估任务上的打分分布与一致性，本附录提供了详细的模型打分热力图。

\begin{figure}[htbp]
    \centering
    \TBDFIG{4 款裁判模型 $\times$ 6 个历史基准 的 CRB-Eval 打分热力图（分别对应 4 个评估维度）}
    \bicaption{不同模型在以前基准上的CRB-Eval分数总结。}{Summary of CRB-Eval scores for different models on previous benchmarks.}
    \label{fig:dataset-heatmeap}
\end{figure}
针对本文纳入审计的每一个历史对话检索基准（包括 QReCC、QuAC、Doc2Dial、TopiOCQA、InSCIT 以及 CORAL），我们分别收集了 4 款参评大语言模型在四维指标上的独立评分。这一系统性的审计过程最终生成了 4 张规格为 $4 \times 6$ 的"裁判模型-评测基准"打分矩阵热力图，如图~\ref{fig:dataset-heatmeap} 所示。

通过对热力图的纵向观察（即同一基准在不同裁判模型下的得分分布）可以发现，尽管 4 款大语言模型在绝对分值上存在微小的固有偏好差异，但它们在不同基准之间给出的相对排序与质量趋势保持了高度的一致性。这从微观层面进一步印证了正文中所采用的“异构多模型集成与逐点评分策略”能够有效稀释单一模型的评估偏见，确保审计结果的客观性与科学性。



% --- 附录 B ---
\section{CORAL 基准审计详细分析}
\label{app:coral_audit}

\textbf{构建机制与局限性。}CORAL 是利用大语言模型自动化构建多轮对话检索基准的早期尝试，其核心逻辑是利用维基百科固有的层级结构（标题--子标题--正文）来生成对话实例。具体而言，该方法利用 GPT-4 将各级标题改写为用户提出的问题，并将标题下的对应正文段落作为系统回答，而金标文档则标注为正文中引用的参考文献。

然而，这种基于静态文档结构的启发式方法实质上是将“人工的局部视野”替换为了“规则的局部视野”。由于缺乏全局性的验证机制，生成的查询往往与标注文档之间存在严重的语义错位，且对话的展开路径直接反映了原始文档的物理组织而非用户的真实意图流转。

通过 \textsc{CRB-Eval} 对 CORAL 原始数据集进行的四维量化审计显示，其数据质量在多个关键维度上显著低于人工标注基准及本文提出的 \textsc{CRB-Bench}。具体评分如下：
\begin{itemize}
    \item \textbf{证据完备性：} \TBD{} / 5.0
    \item \textbf{查询--证据对齐度：} \TBD{} / 5.0
    \item \textbf{答案--证据忠实度：} \TBD{} / 5.0
    \item \textbf{答案质量：} \TBD{} / 5.0
\end{itemize}
其中，查询--证据对齐度与答案--证据忠实度的极低评分反映出该基准中充斥着严重的标注噪声与对齐失效问题。

\textbf{典型对齐失效案例研究。}为了直观展示 CORAL 的质量缺陷，本节选取了数据集中的首个样本进行定性分析。该案例展示了明显的系统性错位：对话主题涉及天文学，而关联的证据文档却是关于数字对象标识符（DOI）系统的技术说明。

\textbf{用户:} What methods are used to predict these geodynamic regimes?  

\textbf{答案:} exoplanets have been directly observed and remotely sensed, [15] but due to their great distance and proximity to obscuring energy sources (the stars they orbit), there is little concrete knowledge of their composition and geodynamic regime. therefore, the majority of information and conjectures made about them come from alternative sources [cite: 

\textbf{金标证据：} doi it uses a managed registry (providing social and technical infrastructure). it does not assume any specific business model for the provision of identifiers or services and enables other existing services to link to it in defined ways...

\textbf{审计结论。}本研究发现，CORAL 生成的查询常与标注的支撑文档存在系统性脱节。在将上述定量与定性证据反馈给 CORAL 团队后，对方承认了原始版本中存在的质量问题，并随后撤回了原始仓库。这一案例深刻表明：脱离全局感知与意图驱动的自动化合成，仅能产出信噪比极低的数据，无法胜任严谨的检索能力评估任务。
% --- 附录 C ---
\section{标注认知边界案例研究}
\label{app:cognitive_boundary}

\textbf{理论概述：认知边界与标注稀疏性}。人工标注在构建多轮对话检索基准时，天然受到“认知边界”（Cognition Boundary）的局限 。标注者在设计评估实例（Evaluation Instances）时，通常采用“以文档为中心”的工作流：即先阅读一篇特定的文档，随后针对该文档的内容构思查询 。

然而，在大规模语料库（Large-scale Corpus）环境下，标注者选定的所谓“金标文档”（Gold Document）往往并非该查询唯一或排他的答案来源 。由于标注者无法在标注过程中全局审视整个语料库，大量同样包含正确答案（甚至更优）的有效证据文档被系统性地遗漏了，这直接导致了标注稀疏性（Annotation Sparsity）问题 。

\textbf{}{典型案例对比分析。}为了直观展示认知边界对数据集质量的影响，表 \ref{tab:sparsity_cases} 对比了各主流基准中人工标注的金标文档与 \textsc{CRB-Eval} 在语料库中识别出的有效替代证据（Alternative Valid Evidence）。

\begin{table}[htbp]
\small
\bicaption{说明标注稀疏性的案例研究}{Case study illustrating annotation sparsity}
\label{tab:sparsity_cases}
\centering
\renewcommand{\arraystretch}{1.5}
\begin{tabularx}{\textwidth}{p{1.6cm} p{3.2cm} X X}
\toprule
\textbf{数据集} & \textbf{查询} & \textbf{金标文档} & \textbf{替代性证据（已检索）} \\
\midrule

\textbf{QReCC} & 
\textit{请告诉我心律不齐的具体类型。} & 
\dots 心律紊乱可能是正常的生理反应 \dots 当肺部无法提取氧气时可能会发生缺氧 \dots 严重贫血 \dots 会降低携氧能力 \dots & 
\dots 心律失常根据其发生部位进行分类 \dots \textbf{心律失常共有十二种类型} \dots 当心房失去协调跳动的能力时，就会发生心房颤动 \dots \\
\midrule

\textbf{Doc2Dial} & 
\textit{我想获取一些关于退休福利计划工具（retirement benefits planner）的信息。} & 
福利计划工具：退休\textbf{在线计算器}（WEP版本）。下方显示的计算器可让您估算您的社会保障福利 \dots 注意：如果您的生日是1月1日 \dots & 
\textbf{福利计划工具} \dots 使用我们的计划工具帮助您更好地了解您的社会保障 \dots \textbf{退休福利：使用我们的退休计划工具来了解}：您如何符合资格 \dots 关于可能的福利 \dots \\
\midrule

\textbf{TopiOCQA} & 
\textit{Victor 什么时候开始出演《年轻和躁动不安》的？} & 
William J. Bell 将 Victor 创作为短期非合约角色，\textbf{于1980年2月8日首演}。Bell 在1997年表示 \dots & 
 Victor Newman 是一个虚构角色 \dots 他\textbf{自1980年起由 Eric Braeden 扮演}。最初作为客串角色，原计划出场八到十二周 \dots \\
\midrule

\textbf{InSCIT} & 
\textit{唯一必须披露痕量过敏原的国家是哪个 \dots？} & 
\dots 尽管如此，目前没有任何标签法律强制要求声明痕量过敏原的存在 \dots \textbf{除了巴西}。 & 
 \textbf{在巴西，自2016年4月起}，当产品未故意添加任何致敏食品时，强制要求声明可能发生交叉污染 \dots \\
\midrule

\textbf{QuAC} & 
\textit{《Death of a Ladies' Man》是一张专辑吗？} & 
Spector 开始重新崭露头角 \dots 制作并参与创作了 Leonard Cohen 在1977年发行的一张备受争议的\textbf{专辑，名为《Death of a Ladies' Man》}。这激怒了许多 Cohen 的忠实粉丝 \dots & 
 所有歌曲均由 Leonard Cohen 创作 \dots 分类：\textbf{1977年专辑}，Leonard Cohen 专辑 \dots Cohen 于1978年出版了《Death of a Lady's Man》一书。它与该专辑没有任何共同之处 \dots \\

\bottomrule
\end{tabularx}
\end{table}

\textbf{案例总结与启示。}通过上述案例可以观察到，语料库中的相关证据往往具有非唯一性：

\textbf{事实冗余性}：在 TopiOCQA 和 InSCIT 中，替代文档包含与金标完全一致的事实信息（如日期“1980”或国家“巴西”），其证明力是相等的。

\textbf{答案直接性}：在 QReCC 和 Doc2Dial 案例中，替代证据有时甚至比原标注文档更能直接、全面地回答用户的特定提问角度。


综上所述，如果评估体系（如 Recall 或 MRR 指标）仅仅依赖由于认知边界产生的单一“金标”，将会错误地惩罚那些具备强大检索能力、能从长尾分布中寻找到最优证据的检索模型，从而导致评估结果的系统性偏差。


\section{CRB-Pipeline工程实现细节}
\label{app:pipeline_details}

本附录详细阐述了 \textsc{CRB-Pipeline} 在实际构建过程中的技术决策，包括多智能体角色的模型对齐实验以及贪心遍历聚类算法的输出示例。

\subsection{模型选择与组合评估}
\label{subsec:app_model_selection}


\begin{figure}[htbp]
    \bicaption{评估由各种模型配对生成的基准的质量，以选择最佳组合}{Evaluating the quality of benchmarks generated by various model pairings to select the best combination}
    \centering
    \TBDFIG{提问者$\times$回答者模型配对热力图（$N\times N$ 交叉矩阵，评分为 CRB-Eval 四维平均）}
    \label{fig:app_model_heatmap}
\end{figure}

在确定对话合成流水线中"提问者（Questioner）"与"回答者（Responder）"的最佳模型组合时，本文对选定的 \TBD{候选模型数} 款开源大语言模型进行了全矩阵交叉评估。实验探索了所有 \TBD{组合总数} 种两两组合可能性，并分别为每对组合生成了包含 \TBD{每对样本数} 个交互轮次的测试数据集 。

这些生成的演示数据集构成了评估不同模型配对效能的实证基础，最终指导我们为特定角色选择表现最优的开源模型。图 \ref{fig:app_model_heatmap} 展示了这些模型组合在四个质量维度（查询相关性、答案质量、标注准确性、响应忠实度）上的打分热力图。



根据评估结果，虽然更强大的模型组合往往能产生更高质量的数据，但本文也观察到部分模型在特定角色（如提问意图的多样性）上展现出独特的优势。这种角色解耦的设计允许研究人员根据计算预算和质量需求灵活调整模型配置。

\subsection{聚类示例}
\label{subsec:app_clustering_example}

为了直观展示贪心遍历聚类算法如何模拟自然的话题演进规律，表 \ref{tab:clustering_example} 提供了一个典型的文档簇示例。该示例展示了围绕“加拿大原住民及其相关主题”构建的语义路径。

在处理过程中，系统首先使用 \texttt{gte-Qwen2-7B} 模型将文本转换为嵌入向量，并利用 \texttt{Faiss} 库构建加权无向图 。随后通过贪心算法寻找 TSP 路径，并按固定步长切分为簇 。如表所示，同一个簇内的文档既保持了高度的语义关联，又在主题上实现了平滑的线性过渡（如从具体的“捕猎路径”逐步演进到宏观的“民族身份定义”）。

\begin{table}[!htbp]
\centering
\bicaption{加拿大原住民及其相关主题的聚类路径示例}{A cluster on Indigenous Peoples and Related Topics in Canada}
\label{tab:clustering_example}
\resizebox{\linewidth}{!}{
\begin{tabular}{lp{10cm}}
\toprule
\textbf{文档标题} & \textbf{摘要内容（精简）} \\ \midrule
Trapline & 介绍捕猎路径的传统知识与文化意义，以及加拿大各省对其正式登记的管理制度。 \\ \midrule
Economy of Saskatoon & 描述萨斯卡通市的“城市保留地合作模式”，其中原住民土地由城市提供服务并共同纳税。 \\ \midrule
Indian Reserve & 定义了加拿大《印第安人法案》下的保留地概念，并说明其与原住民土地所有权的差异。 \\ \midrule
First Nations in Canada & 详细说明“第一民族”这一称谓的演变过程及其法律地位，涵盖 600 多个公认的政府/部落。 \\ \midrule
Indigenous peoples & 统称第一民族、因纽特人和梅蒂人，并讨论了宪法法律术语的演进趋势。 \\ \midrule
Métis & 专门探讨梅蒂人的独特性，包括其作为欧洲定居者与原住民后裔的混合文化身份 。 \\ \midrule
Cree & 介绍克里族原住民的历史，以及梅蒂人身份中经常包含的克里族血统背景 。 \\ \bottomrule
\end{tabular}
}
\end{table}

通过这种方式，\textsc{CRB-Pipeline} 成功克服了传统聚类方法产生的文档重叠与跨簇冗余问题，为后续多轮对话生成提供了纯净且具备逻辑深度的上下文基底。


\section{各角色Prompt模板}
\label{app:prompts}

本附录汇集了本文所提框架中涉及的所有关键 Prompt 模板。

\textsc{CRB-Pipeline} 数据合成模板。用于指导多智能体协作生成多轮对话检索数据的核心提示词。如表\ref{tab:prompt1}到\ref{tab:polisher}所示。

\textbf{\textsc{CRB-Eval} 审计裁判模板。}用于自动化量化评估基准质量的细粒度打分提示词模板。如表\ref{tab:question_relevance_prompt}到\ref{tab:Answer Quality Prompt}所示。
\begin{table}[!htbp]
\bicaption{提问者提示词}{Questioner Prompt}
\small
\input{assets/prompt/QUERY_WO_HISTORY}

    \label{tab:prompt1}
\end{table}
\begin{table}
\bicaption{回答者提示词}{Reasoner Prompt}
\small
\input{assets/prompt/RESPONSE}
    \label{tab:prompt2}
    
\end{table}


\begin{table}[!htbp]
\bicaption{润色者提示词}{Polisher Prompt}
\small
\input{assets/prompt/PHRASE}
    \label{tab:polisher}
\end{table}

\begin{table}[!htbp]
\bicaption{查询-证据对齐度提示词}{Query-Evidence Alignment Prompt}
\small
\input{assets/prompt/JUDGE_QUERYER_TAG}
    \label{tab:question_relevance_prompt}
\end{table}

\begin{table}[!htbp]
\bicaption{证据完备性提示词}{Evidence Completeness Prompt}
\small
\input{assets/prompt/JUDGE_QUERY_RELATED}
    \label{tab:Annotation Correctness Prompt}
\end{table}



\begin{table}[!htbp]
\bicaption{答案-证据忠实度提示词}{Answer-Evidence Faithfulness Prompt}
\small
\input{assets/prompt/JUDGE_RESPONSE_FAITHFUL}
    \label{tab:RESPONSE_FAITHFUL}
\end{table}



\begin{table}[!htbp]
\small
\bicaption{答案质量}{Answer Quality Prompts}
\input{assets/prompt/JUDGE_ANSWER_QUALITY}
    \label{tab:Answer Quality Prompt}
\end{table}



\section{合成对话示例}
\label{app:synthetic_dialogue}

本附录展示了由 \textsc{CRB-Pipeline} 全自动合成的一段完整的多轮对话检索示例。该示例生动呈现了提问者智能体（Questioner）的话题转移、指代消解与省略现象，以及回答者智能体（Responder）基于长程上下文的详尽推理（Chain-of-Thought）与事实锚定能力。

用CRB-Pipeline生成的对话如下所示（‘U’和‘A’分别代表用户和智能体）。该对话共包含 8 个交互轮次，涵盖了线性话题深化、硬话题跳跃、以及意图回溯等复杂的真实认知行为。

U: 谁在莱比锡大学建立了第一个心理学实验室？

A: 莱比锡大学的第一个心理学实验室由威廉·冯特（Wilhelm Wundt）建立。冯特是一位受过训练的生理学家，他对感官知觉进行了研究，并在化学进展的启发下致力于研究心理过程的基本组成部分。他的实验室将实验心理学带向了世界。

U: 那么“心理物理学”（psychophysics）这个词是谁首创的？
A: “心理物理学”一词是由古斯塔夫·西奥多·费希纳（Gustav Theodor Fechner）首创的。

U: 等等，刚才说在莱比锡大学建立那个实验室的人是谁来着？

A: 莱比锡大学的心理学实验室是由威廉·冯特建立的。冯特深受古斯塔夫·费希纳和赫尔曼·冯·亥姆霍兹在感官知觉方面工作的影响。在化学进展的启发下，冯特致力于使用实验方法研究心理过程的基本组成部分。他建立于1879年的实验室被广泛认为是第一个专门从事实验心理学研究的正式实验室，标志着心理学作为一门科学学科历史上的一个重要里程碑。

U: 噢，顺便问一下，雨果·明斯特伯格（Hugo Münsterberg）出生在哪里？

A: 雨果·明斯特伯格出生在但泽（Danzig），即现在的波兰格但斯克。在他出生时，但泽是西普鲁士的一个港口城市。

U: 奥特拉皮兹（Otter Rapids）社区的人口是多少？

A: 奥特拉皮兹社区的人口超过了 1,000 人。

U: 那里有什么样的生活设施？

A: 奥特拉皮兹社区设施齐全，能够满足居民的各种需求。这些设施包括银行和邮局等基本服务。在餐饮和购物方面，有一家小吃店和一家杂货店。娱乐设施也非常丰富，居民可以享受台球厅、保龄球馆、冰壶场和冰球场。安全和健康服务由消防站和一家拥有 10 张床位的医院提供。此外，社区还有一个礼堂，可能用于举办各种活动和集会。这些设施确保了这个拥有 1000 多人的社区能够获得必要的服务和休闲活动。

U: 那么，科托溪（Coteau Creek）水电站是什么时候开始运营的？

A: 为了确定科托溪水电站开始运营的年份，我们可以按照以下步骤进行推理：1. 确认建设时间线：为科托溪水电站供水的加德纳大坝（Gardiner Dam）建于1959年至1967年之间。2. 发电站的建设：发电站本身的建设需要近三年的时间。3. 运营年份：鉴于发电站的建设是在1967年加德纳大坝完工后开始的，并且耗时近三年，我们可以通过在1967年的基础上增加三年得出其运营年份。因此，科托溪水电站于1969年开始运营。

U: 那么奥特拉皮兹的第一台发电机组是什么时候并网发电的？

A: 为了确定奥特拉皮兹发电站第一台发电机组并网的年份，我们可以按照以下步骤进行：1. 确认开发时间线：奥特拉皮兹发电站的开发经历了几个关键的里程碑。初步的调查在1945年进行，1951年又收集了更多的现场信息。2. 批准与建设：该项目于1958年10月获得了相关法规的批准。这项批准为发电机组的实际建设和安装铺平了道路。3. 首台机组并网：奥特拉皮兹发电站的第一台发电机组于1961年9月26日并网。因此，第一台发电机组是在1961年并网发电的。



% --- 附录结束 ---
% --- 附录结束 ---

\cleardoublepage[plain]



\chapter[致谢]{致\quad 谢}\chaptermark{致\quad 谢}% syntax: \chapter[目录]{标题}\chaptermark{页眉}
% \thispagestyle{noheaderstyle}% 如果需要移除当前页的页眉
%\pagestyle{noheaderstyle}% 如果需要移除整章的页眉



\reviewORprint{
}{
   
}

\chapter{攻硕期间参与项目、发表论文及获奖情况}


\section*{专利:}
\begin{enumerate}[leftmargin=*]
    \item 一种基于检索增强生成的语音大模型问答
方法\quad 三作
\end{enumerate}

\reviewORprint{
    \section*{第一作者发表/录用学术论文:}
    \begin{enumerate}[leftmargin=*]
        \item Title[J]Journal, Year, Volume(Number):Pages. (\textbf{SCI, JCR1, 论文第三章})
    \end{enumerate}
    \section*{第一作者在审学术论文:}
    \begin{enumerate}[leftmargin=*]
        \item Title[J]Journal, Year, Volume(Number):Pages. (\textbf{SCI, JCR1, 论文第三章})
    \end{enumerate}
    \section*{通讯作者发表/录用学术论文:}
    \begin{enumerate}[leftmargin=*]
        \item Title[J]Journal, Year, Volume(Number):Pages. (\textbf{SCI, JCR1, 第二作者/通讯作者})
    \end{enumerate}
    \section*{合作作者发表/录用学术论文:}
    \begin{enumerate}[leftmargin=*]
        \item Title[J]Journal, Year, Volume(Number):Pages. (\textbf{SCI, JCR1, 第二作者})
    \end{enumerate}
}{
    \section*{学术论文:}
    \begin{enumerate}[leftmargin=*]
        \item MTR-Suite: A Framework for Evaluating and Synthesizing Conversational Retrieval Benchmarks \quad ACL 2026 main\quad 第二作者
        \item Enhancing Neural Machine Translation Through Target Language Data: A kNN-LM Approach for Domain Adaptation \quad ACL 2025 main\quad 第四作者
    \end{enumerate}
}

\section*{获奖情况:}
\begin{enumerate}[leftmargin=*]
    \item 获得2023-2024年度东北大学硕士研究生一等奖学金
    \item 获得2025-2026年度东北大学硕士研究生二等奖学金
\end{enumerate}

\cleardoublepage[plain]% 让文档总是结束于偶数页，可根据需要设定页眉页脚样式，如 [noheaderstyle]

% other information
% \showthe\artxfontset
\end{sloppypar}
\end{document}
%---------------------------------------------------------------------------%


% \usepackage[most]{tcolorbox}